\documentclass[11pt]{article}

\usepackage{amsmath,amssymb,amsthm,amsfonts,mathtools,bm,enumitem,longtable,array}
% Use the standard algorithm packages when available. The fallbacks keep this
% source compilable in minimal TeX installations without changing the algorithms.
\IfFileExists{algorithm.sty}{\usepackage{algorithm}}{%
  \usepackage{float}
  \floatstyle{ruled}
  \newfloat{algorithm}{htbp}{loa}
  \floatname{algorithm}{Algorithm}
}
\IfFileExists{algpseudocode.sty}{\usepackage{algpseudocode}}{%
  \newcounter{ALGline}
  \newlength{\ALGindent}
  \newenvironment{algorithmic}[1][1]{%
    \setlength{\ALGindent}{0pt}%
    \begin{list}{\scriptsize\arabic{ALGline}:}{%
      \usecounter{ALGline}\setlength{\leftmargin}{2em}%
      \setlength{\labelwidth}{1.5em}\setlength{\labelsep}{.5em}%
      \setlength{\itemsep}{0pt}\setlength{\parsep}{0pt}}}{\end{list}}
  \newcommand{\State}{\item\hspace*{\ALGindent}}
  \newcommand{\Require}{\item[]\textbf{Require:} }
  \newcommand{\Return}{\textbf{return}}
  \newcommand{\If}[1]{\State\textbf{if} ##1 \textbf{then}\addtolength{\ALGindent}{1em}}
  \newcommand{\EndIf}{\addtolength{\ALGindent}{-1em}\State\textbf{end if}}
  \newcommand{\For}[1]{\State\textbf{for} ##1 \textbf{do}\addtolength{\ALGindent}{1em}}
  \newcommand{\EndFor}{\addtolength{\ALGindent}{-1em}\State\textbf{end for}}
  \newcommand{\Repeat}{\State\textbf{repeat}\addtolength{\ALGindent}{1em}}
  \newcommand{\Until}[1]{\addtolength{\ALGindent}{-1em}\State\textbf{until} ##1}
}
\usepackage[top=1in,bottom=1in,left=1in,right=1in]{geometry}
\usepackage{titletoc}
\usepackage{setspace}
\usepackage[colorlinks,allcolors=blue,linktoc=all,bookmarksopen=true,bookmarksnumbered=true]{hyperref}
\usepackage[round]{natbib}
\usepackage[nameinlink,noabbrev]{cleveref}
\usepackage{microtype}
\onehalfspacing
\allowdisplaybreaks


\titlecontents{section}[0em]{\vspace{7.5pt}\bfseries}{\thecontentslabel\quad}{}{\titlerule*[0.75pc]{-}\contentspage}
\titlecontents{subsection}[2em]{}{\thecontentslabel\quad}{}{\titlerule*[0.75pc]{-}\contentspage}
\titlecontents{subsubsection}[3em]{}{\thecontentslabel\quad}{}{\titlerule*[0.75pc]{-}\contentspage}

% Theorem environments.
\newtheorem{theorem}{Theorem}[section]
\newtheorem{definition}{Definition}[section]
\newtheorem{proposition}{Proposition}[section]
\newtheorem{lemma}{Lemma}[section]
\newtheorem{corollary}{Corollary}[section]
\newtheorem{remark}{Remark}[section]
\newtheorem{assumption}{Assumption}

% Hyperlinked theorem-like references. With `nameinlink`, the entire
% reference (e.g. ``Proposition 5.1'') is clickable in the PDF.
\crefname{theorem}{Theorem}{Theorems}
\Crefname{theorem}{Theorem}{Theorems}
\crefname{lemma}{Lemma}{Lemmas}
\Crefname{lemma}{Lemma}{Lemmas}
\crefname{proposition}{Proposition}{Propositions}
\Crefname{proposition}{Proposition}{Propositions}
\crefname{corollary}{Corollary}{Corollaries}
\Crefname{corollary}{Corollary}{Corollaries}
\crefname{definition}{Definition}{Definitions}
\Crefname{definition}{Definition}{Definitions}
\crefname{assumption}{Assumption}{Assumptions}
\Crefname{assumption}{Assumption}{Assumptions}

% Convenience macros.
\newcommand{\R}{\mathbb{R}}
\newcommand{\E}{\mathbb{E}}
\newcommand{\Pp}{\mathbb{P}}
\newcommand{\F}{\mathcal{F}}
\newcommand{\X}{\mathcal{X}}
\newcommand{\Ncal}{\mathcal{N}}
\newcommand{\eps}{\varepsilon}
\newcommand{\op}{\mathrm{op}}
\newcommand{\HS}{\mathrm{HS}}
\newcommand{\tr}{\operatorname{tr}}
\newcommand{\ran}{\operatorname{ran}}
\newcommand{\supp}{\operatorname{supp}}
\newcommand{\argmin}{\operatorname*{arg\,min}}
\newcommand{\argmax}{\operatorname*{arg\,max}}
\newcommand{\esssup}{\operatorname*{ess\,sup}}
\newcommand{\inner}[2]{\left\langle #1,#2\right\rangle}
\newcommand{\norm}[1]{\left\lVert #1\right\rVert}
\newcommand{\abs}[1]{\left\lvert #1\right\rvert}
\newcommand{\Kx}[1]{K_{#1}}
\newcommand{\whV}{\widehat V_n}
\newcommand{\whw}{\widehat\omega}
\newcommand{\whf}{\widehat f}
\newcommand{\whth}{\widehat\theta_{\mathrm{GM}}}
\newcommand{\preceqL}{\preceq}
\newcommand{\succeqL}{\succeq}

\begin{document}

\title{\textbf{Debiased Inference in RKHS}\\[4pt]\large Extension to Latent Gaussian Confounding}
\author{Samuel N.E. Lee\thanks{Department of Statistics and Data Science, National University of Singapore.} \and Alexander Giessing\thanks{Department of Statistics and Data Science, National University of Singapore.}}
\date{September 14, 2026}
\maketitle
\tableofcontents
\newpage

% ============================================================
\section{Notation}
% ============================================================

\subsection*{Model and RKHS quantities}
\begin{itemize}[leftmargin=2em]
  \item $(X_i,Y_i,E_i)_{i=1}^n$: observations in the second fold $I_2$, with $X_i\in\X=\R^d$, $Y_i\in\R$, and $E_i\in\{e,f\}$. The first fold $I_1$ has size $m$, with $m/n\to\rho\in(0,\infty)$.
  \item $\nu_g:=\mathcal L(X\mid E=g)$ and $\nu:=\pi_e\nu_e+\pi_f\nu_f$: the pooled marginal law of $X$, where $\pi_g>0$ and $\pi_e+\pi_f=1$; $L_2(\nu)$ is the associated square-integrable function space.
  \item $\F$: a separable reproducing kernel Hilbert space (RKHS) on $\X$ with continuous kernel $K$.
  \item $K_x:=K(\cdot,x)\in\F$: the kernel section at $x$, so that $f(x)=\inner{f}{K_x}_{\F}$.
  \item $f_0\in\F$: the true regression function; $Y_i=f_0(X_i)+H_i^\top\delta_0+\eps_i$.
  \item $x_0\in\X$: a fixed query point; $\theta_0:=f_0(x_0)$ is the target functional.
  \item $\bm Y=(Y_1,\ldots,Y_n)^\top$, $\widetilde{\bm Y}=(Y_i-\widehat m_{E_i}(X_i))_{i=1}^n$, and $\bm U=(U_1,\ldots,U_n)^\top$.
  \item $K\in\R^{n\times n}$: the empirical Gram matrix, $K_{ij}:=K(X_i,X_j)$.
  \item $k_0\in\R^n$: $(k_0)_i:=K(x_0,X_i)$; $k_{00}:=K(x_0,x_0)$.
  \item For $f\in\F$, $f(\bm X):=(f(X_1),\ldots,f(X_n))^\top$.
\end{itemize}

\subsection*{Confounding and environment quantities}
\begin{itemize}[leftmargin=2em]
  \item $\Omega_g:=\Gamma^\top\Sigma_H\Gamma+\Sigma_g$, $P_g:=\Omega_g^{-1}$, and $c:=\Gamma^\top\Sigma_H\delta_0$.
  \item $w_g(x):=P_g(x-\mu_g)$ and $m_g(x):=c^\top P_g(x-\mu_g)$: the conditional mean contribution of the latent confounder in environment $g$.
  \item $U_i:=H_i^\top\delta_0+\eps_i-m_{E_i}(X_i)$: the corrected error, satisfying $\E[U_i\mid X_i,E_i]=0$.
  \item $\sigma_{U,g}^2:=\operatorname{Var}(U_i\mid X_i,E_i=g)>0$ and $t_g:=\sigma_{U,g}^{-2}$.
  \item $\widehat c$, $\widehat P_g$, $\widehat\mu_g$, and $\widehat\sigma_{U,g}^2$: estimates constructed using only $I_1$; $\widehat t_g:=(\widehat\sigma_{U,g}^2)^{-1}$.
  \item $q_g(x):=\Pr(E=g\mid X=x)$: the posterior environment probability in the working mixture.
\end{itemize}

\subsection*{Operators and spectral quantities}
\begin{itemize}[leftmargin=2em]
  \item $I_\nu:\F\to L_2(\nu)$: canonical embedding, $I_\nu f=[f]_\nu$; $S_\nu:=I_\nu^*$.
  \item $T_\nu:=I_\nu S_\nu:L_2(\nu)\to L_2(\nu)$ and $C_\nu:=S_\nu I_\nu:\F\to\F$.
  \item $C_g:=\E[K_X\otimes K_X\mid E=g]$ and $C_\nu=\pi_eC_e+\pi_fC_f$.
  \item $C_{\nu,n}:=n^{-1}\sum_{i=1}^n K_{X_i}\otimes K_{X_i}$: empirical covariance operator.
  \item $D_\nu:=\pi_et_eC_e+\pi_ft_fC_f$ and $D_{\nu,n}:=n^{-1}\sum_{i=1}^n\widehat t_{E_i}K_{X_i}\otimes K_{X_i}$: population and empirical precision-weighted covariance operators.
  \item $\Ncal_{D,x}(\eta):=\inner{K_x}{(D_\nu+\eta I)^{-1}K_x}_{\F}$ and $\widehat\Ncal_{D,x}(\eta):=\inner{K_x}{(D_{\nu,n}+\eta I)^{-1}K_x}_{\F}$.
  \item $W:=\operatorname{diag}(\widehat\sigma_{U,E_1}^2,\ldots,\widehat\sigma_{U,E_n}^2)$ and $\widetilde S_n:=W^{-1/2}S_n$.
  \item $\rho_j$ and $\psi_j$: the positive eigenvalues and orthonormal eigenvectors of $D_\nu$; if its rank is finite, $\rho_{\min}^+>0$ denotes its smallest positive eigenvalue.
  \item $r_\nu:=\operatorname{rank}(C_\nu)\in\mathbb N\cup\{\infty\}$: the number of non-zero eigenvalues of $C_\nu$.
  \item $\mathcal J_\nu:=\{1,\ldots,r_\nu\}$ if $r_\nu<\infty$, and $\mathcal J_\nu:=\mathbb N$ if $r_\nu=\infty$.
  \item $\mu_1\ge\mu_2\ge\cdots>0$: the non-zero eigenvalues shared by $T_\nu$ and $C_\nu$, indexed by $j\in\mathcal J_\nu$.
  \item $\F_\nu:=(\ker I_\nu)^\perp=\overline{\ran(C_\nu)}$: the identifiable subspace of $\F$; when $r_\nu<\infty$, $\F_\nu=\ran(C_\nu)$.
  \item $\Ncal_\nu(\eta):=\tr\{C_\nu(C_\nu+\eta I)^{-1}\}$: global effective dimension.
  \item $\Ncal_x(\eta):=\inner{K_x}{(C_\nu+\eta I)^{-1}K_x}_{\F}$: population pointwise regularized leverage score.
  \item $\widehat{\Ncal}_x(\eta):=\inner{K_x}{(C_{\nu,n}+\eta I)^{-1}K_x}_{\F}$: empirical regularized leverage score.
  \item For bounded self-adjoint operators $A,B$, $A\preceqL B$ means $B-A$ is positive semidefinite.
  \item $\norm{A}_{\op}$, $\norm{A}_{\HS}$, and $\norm{A}_{1}$ denote the operator, Hilbert--Schmidt, and trace norms, respectively.
\end{itemize}

\subsection*{Regularization and debiasing parameters}
\begin{itemize}[leftmargin=2em]
  \item $\lambda>0$: KRR penalty 
  % $\lambda_n\asymp n^{-1/(\beta+p)}$ in the minimax-rate theorem.
  \item $\eta>0$: regularization parameter for the inverse problem defining the debiasing weights; may be written as $\eta_n$ to emphasize dependence on $n$; $\eta_n\asymp (\log n)^2/n$ in the infinite-rank case and $\eta_n\asymp n^{-1}$ in the finite-rank case in the inferential theory.
  \item $\gamma\ge0$: RKHS alignment tolerance in the primal GM program; may be written as $\gamma_n$ to emphasize dependence on n.
  \item $\mu\ge0$: scalar Lagrange multiplier in the dual GM program.
  \item $p\in(0,1)$: global eigenvalue/effective-dimension exponent in the infinite-rank case $r_\nu=\infty$.
  \item $\beta\in(1,2]$: source-smoothness exponent for $f_0$ in the infinite-rank case. 
  \item $a\in(0,1)$: local regular-variation exponent at the fixed query point $x_0$ in the infinite-rank case. In the finite-rank case, the analogous leverage is bounded and behaves as the case $a=0$.
  \item $\Delta_\omega:=K_{x_0}-n^{-1}\sum_{i=1}^n\omega_iK_{X_i}$: alignment residual generated by weights $\omega\in\R^n$.
  \item $\widehat V_n:=n^{-1}\sum_{i=1}^n\widehat\sigma_{U,E_i}^2\widehat\omega_i^2$: observable variance proxy; $V_n$ denotes its population version of the variance proxy.
\end{itemize}

\subsection*{Asymptotic notation}
For nonnegative sequences $A_n,B_n$, $A_n=O(B_n)$ means $A_n/B_n$ is eventually bounded, $A_n\asymp B_n$ means both $A_n=O(B_n)$ and $B_n=O(A_n)$, and $A_n=o(B_n)$ means $A_n/B_n\to0$. We write $A_n=O_p(B_n)$ if $A_n/B_n$ is bounded in probability and $A_n=o_p(B_n)$ if $A_n/B_n\to0$ in probability. The notation $\overset{d}{\to}$ denotes convergence in distribution and $\overset{p}{\to}$ convergence in probability. Constants denoted by $C,c,A,D$ may change between displays and are independent of $n$.

% ============================================================
\section{Definitions}
% ============================================================

\begin{definition}[RKHS and rank-one operators]
A Hilbert space $\F$ of real-valued functions on $\X$ is an RKHS with kernel $K$ if $K_x\in\F$ for every $x\in\X$ and
\[
  f(x)=\inner{f}{K_x}_{\F},\qquad f\in\F.
\]
For $g,h\in\F$, the rank-one operator $g\otimes h:\F\to\F$ is
\[
  (g\otimes h)f:=\inner{f}{h}_{\F}g.
\]
In particular, $K_x\otimes K_x$ is a positive semidefinite rank-one operator in the space of Hilbert--Schmidt operators.
\end{definition}

\begin{definition}[Canonical embedding and covariance operators]
Let $I_\nu:\F\to L_2(\nu)$ map $f$ to its $\nu$-equivalence class $[f]_\nu$, and let $S_\nu:=I_\nu^*$. Define
\[
  T_\nu:=I_\nu S_\nu:L_2(\nu)\to L_2(\nu),\qquad
  C_\nu:=S_\nu I_\nu:\F\to\F.
\]
Equivalently,
\[
  C_\nu f=\E[f(X)K_X]=\E[(K_X\otimes K_X)f].
\]
Under a bounded measurable kernel, these operators are positive, self-adjoint, and trace class, hence compact, and they have the same non-zero eigenvalues. We write
\[
  r_\nu:=\operatorname{rank}(C_\nu)\in\mathbb N\cup\{\infty\},
  \qquad
  \mathcal J_\nu
  :=
  \begin{cases}
    \{1,\ldots,r_\nu\}, & r_\nu<\infty,\\
    \mathbb N, & r_\nu=\infty.
  \end{cases}
\]
The identifiable subspace is
\[
  \F_\nu:=(\ker I_\nu)^\perp=\overline{\ran(C_\nu)}.
\]
Since $\F$ is separable, so is the closed subspace $\F_\nu$. Moreover, $K_X\in\F_\nu$ almost surely. This is because for a countable orthonormal basis $(u_j)$ of $\ker I_\nu$, $\inner{K_X}{u_j}_{\F}=u_j(X)=0$ almost surely for every $j$.
If $r_\nu<\infty$, then $\F_\nu=\ran(C_\nu)$ is finite dimensional. 
\end{definition}

Let $(\mu_j)_{j\in\mathcal J_\nu}$ denote the non-zero eigenvalues, ordered as $\mu_1\ge\mu_2\ge\cdots>0$. We use functions $(e_j)_{j\in\mathcal J_\nu}\subset\F$ such that $([e_j]_\nu)_{j\in\mathcal J_\nu}$ is an orthonormal basis of $\overline{\ran(I_\nu)}\subseteq L_2(\nu)$ and $(\mu_j^{1/2}e_j)_{j\in\mathcal J_\nu}$ is an orthonormal basis of $\F_\nu$. For later spectral calculations, write $\phi_j:=\mu_j^{1/2}e_j$, so that $(\phi_j)_{j\in\mathcal J_\nu}$ is an orthonormal basis of $\F_\nu$ and $C_\nu\phi_j=\mu_j\phi_j$.

If $r_\nu=\infty$, then by the spectral theorem,
\begin{align}
  T_\nu
  &=\sum_{j\ge1}\mu_j\inner{\,\cdot\,}{[e_j]_\nu}_{L_2(\nu)}[e_j]_\nu,\\
  C_\nu
  &=\sum_{j\ge1}\mu_j\inner{\,\cdot\,}{\mu_j^{1/2}e_j}_{\F}\mu_j^{1/2}e_j.
\end{align}
If $r_\nu<\infty$, the same spectral representations hold with the finite sums
\begin{align}
  T_\nu
  &=\sum_{j=1}^{r_\nu}\mu_j\inner{\,\cdot\,}{[e_j]_\nu}_{L_2(\nu)}[e_j]_\nu,\\
  C_\nu
  &=\sum_{j=1}^{r_\nu}\mu_j\inner{\,\cdot\,}{\mu_j^{1/2}e_j}_{\F}\mu_j^{1/2}e_j.
\end{align}


\begin{definition}[Power spaces]
For $\alpha\ge0$, define
\[
  [\F]^\alpha_\nu
  :=\left\{
      \sum_{j\in\mathcal J_\nu}a_j\mu_j^{\alpha/2}[e_j]_\nu:
      (a_j)_{j\in\mathcal J_\nu}\in\ell_2(\mathcal J_\nu)
    \right\},
\]
with $\alpha$-power norm:
\[
  \left\|\sum_{j\in\mathcal J_\nu}a_j\mu_j^{\alpha/2}[e_j]_\nu\right\|_{[\F]^\alpha_\nu}
  :=\left(\sum_{j\in\mathcal J_\nu}a_j^2\right)^{1/2}.
\]
For $0<\eta<\alpha$, $[\F]^\alpha_\nu\hookrightarrow[\F]^\eta_\nu\hookrightarrow L_2(\nu)$, and $[\F]^\alpha_\nu=\ran(T_\nu^{\alpha/2})$ on the identifiable subspace.
\end{definition}

\begin{definition}[Effective dimensions and leverage scores]
For $\eta>0$, the global effective dimension is
\[
  \Ncal_\nu(\eta):=\tr\!\left(C_\nu(C_\nu+\eta I)^{-1}\right)
  =\sum_{j\in\mathcal J_\nu}\frac{\mu_j}{\mu_j+\eta}.
\]


For a point $x\in\X$, the population and empirical regularized leverage scores are
\[
  \Ncal_x(\eta):=\inner{K_x}{(C_\nu+\eta I)^{-1}K_x}_{\F},\qquad
  \widehat\Ncal_x(\eta):=\inner{K_x}{(C_{\nu,n}+\eta I)^{-1}K_x}_{\F}.
\]
The first quantity can be treated as the local analogue of the effective dimension at $x$ due to the identity $\Ncal_x(\eta) = \tr((K_x \otimes K_x)(C_\nu + \eta I)^{-1})$.
\end{definition}

\begin{definition}[Sampling operator and empirical covariance]
\label{defn:sampling_operator}
Define
\[
  S_n:\F\to\R^n,\qquad
  S_nf:= (f(X_1),\ldots,f(X_n))^\top.
\]
Then, on $\R^n$ equipped with the inner product $\inner{u}{v}_n := \frac{1}{n}u^\top v$, 
\[
  S_n^*u=\frac1{n}\sum_{i=1}^n u_iK_{X_i},\qquad
  C_{\nu,n}=S_n^*S_n,\qquad
  S_nS_n^*=\frac1nK.
\]
The precision-weighted sampling operator satisfies
\[
 \widetilde S_n^*\widetilde S_n=D_{\nu,n},\qquad
 \widetilde S_n\widetilde S_n^*=\frac1nW^{-1/2}KW^{-1/2}.
\]
\end{definition}

\begin{definition}[Kernel ridge regression]
For $\lambda>0$, the KRR estimator based on the corrected responses $\widetilde Y_i:=Y_i-\widehat m_{E_i}(X_i)$ is
\begin{equation}
  \widehat f_\lambda
  :=\argmin_{f\in\F}
  \left\{
    \frac1n\sum_{i=1}^n(\widetilde Y_i-f(X_i))^2+\lambda\norm{f}_{\F}^2
  \right\}.
  \label{eq:krr_definition}
\end{equation}
By the representer theorem \cite{wahba1990spline,kimeldorf1971some},
\[
  \widehat f_\lambda(\cdot)=\sum_{i=1}^n\widehat\alpha_iK(X_i,\cdot),\qquad
  \widehat\alpha=(K+n\lambda I_n)^{-1}\widetilde{\bm Y}.
\]
\end{definition}

\begin{definition}[Alignment residual and minimum feasible tolerance]
For $\omega\in\R^n$, define
\[
  \Delta_\omega
  :=K_{x_0}-\frac1n\sum_{i=1}^n\omega_iK_{X_i}.
\]
The minimum achievable tolerance is
\[
  \gamma_{\min,n}:=\inf_{\omega\in\R^n}\norm{\Delta_\omega}_{\F}.
\]
Writing the constraint in Gram-matrix form gives
\begin{equation}
  \norm{\Delta_\omega}_{\F}^2
  =k_{00}-\frac2n k_0^\top\omega+\frac1{n^2}\omega^\top K\omega.
  \label{eq:alignment_gram_form}
\end{equation}
If $K$ is singular, let $K^\dagger$ denote its Moore--Penrose pseudoinverse. Since the augmented Gram matrix is positive semidefinite, $k_0\in\ran(K)$, and therefore
\begin{equation}
  \gamma_{\min,n}^2=k_{00}-k_0^\top K^\dagger k_0.
  \label{eq:gamma_min_closed}
\end{equation}
When $K$ is invertible, $K^\dagger=K^{-1}$.
\end{definition}

% ============================================================
\section{Model and Assumptions}
% ============================================================

\begin{assumption}[Generating model and RKHS structure]\label{assump:model}
The observations $(X_i,Y_i,E_i)$ are independent across $i$ and satisfy
\begin{align}
 Y_i&=f_0(X_i)+H_i^\top\delta_0+\eps_i,\label{eq:model}\\
 X_i\mid(E_i=g)&=\Gamma^\top H_i+\xi_{i,g}+\mu_g,\qquad g\in\{e,f\}.
 \label{eq:confounding_covariates}
\end{align}
Here $H_i\sim N(0,\Sigma_H)$, $\xi_{i,g}\sim N(0,\Sigma_g)$, and $H_i$ and $\xi_{i,g}$ are independent conditional on the environment. The environment is independent of $(H_i,\eps_i)$, and $\eps_i$ is a mean-zero sub-exponential random variable independent of $(H_i,\xi_{i,E_i},E_i)$. The dimensions are fixed and
\[
 \Omega_g:=\Gamma^\top\Sigma_H\Gamma+\Sigma_g\succ0,
 \qquad P_g:=\Omega_g^{-1},\qquad c:=\Gamma^\top\Sigma_H\delta_0.
\]
The observations are sampled independently within the two environments with fixed proportions $\pi_g$.

% We randomly split the observations into two independent folds $I_1$ and $I_2$ with sizes $m$ and $n$, where $m/n\to\rho\in(0,\infty)$, while keeping the environment proportions the same in each fold so $n_g/n=m_g/m=\pi_g$.  Nuisance quantities are estimated using only $I_1$, and the GM estimator is constructed using $I_2$. Within $I_2$, we relabel the observations as $i=1,\ldots,n$.

The true regression function satisfies $f_0\in\F$, where $\F$ is a separable RKHS. The input space is $\X=\R^d$, the marginal distribution $\nu=\pi_e\nu_e+\pi_f\nu_f$ has full support on $\X$, and $K$ is continuous and uniformly bounded:
\begin{equation}
 \sup_{x\in\X}K(x,x)\leq\kappa^2<\infty.
 \label{eq:bounded_kernel}
\end{equation}
% Under these conditions, the canonical embedding is injective and
% \[
%  \F=\F_\nu=(\ker I_\nu)^\perp=\overline{\ran(C_\nu)}.
% \]

\end{assumption}

\begin{remark}
Continuity of $K$ and full support of $\nu$ make the canonical embedding $I_\nu$ injective. Hence $\F_\nu=\F$.
\end{remark}

\begin{remark}
By Gaussian conditioning,
\begin{equation}
 \E[Y\mid X=x,E=g]=f_0(x)+c^\top P_g(x-\mu_g).
 \label{eq:conditional_mean_confounding}
\end{equation}
Define
\begin{equation}
 U_i:=H_i^\top\delta_0+\eps_i-c^\top P_{E_i}(X_i-\mu_{E_i}).
 \label{eq:corrected_error}
\end{equation}
Then
\[
 Y_i-c^\top P_{E_i}(X_i-\mu_{E_i})=f_0(X_i)+U_i,
 \qquad \E[U_i\mid X_i,E_i]=0.
\]
The conditional variance is constant within each environment:
\begin{equation}
 \sigma_{U,g}^2
 =\operatorname{Var}(\eps_i)+\delta_0^\top\Sigma_H\delta_0-c^\top P_gc.
 \label{eq:environment_variance}
\end{equation}
Moreover, the residual $U$ is independent of $X$ within each environment. Since $\eps_i$ is a centered sub-exponential random variable, $U_i\mid(X_i,E_i=g)$ is also a centered sub-exponential random variable.

\end{remark}

\begin{assumption} 
[Full rank of $\Sigma_e - \Sigma_f$]\label{assump:full_rank_of_cov_mat_diff}
The lowest eigenvalue of $\Sigma_e - \Sigma_f$ is bounded away from zero
\begin{equation}
 \lambda_{\min}(\Sigma_e - \Sigma_f) > 0
 \label{eq:full_rank_of_cov_mat_diff}
\end{equation}
\end{assumption}

\begin{remark}
For two environments, if $q_g(x):=\Pr(E=g\mid X=x)$, then
\[
 M=\E\!\left[q_e(X)q_f(X)\{(P_e-P_f)X -(P_e\mu_e - P_f\mu_f) \}\{(P_e-P_f)X -(P_e\mu_e - P_f\mu_f)\}^\top\right].
\]
\noindent
For every non-zero $u \in \R^d$,

\[u^\top M u = \E[q_e(X)q_f(X)\{u^\top\left((P_e-P_f)X -(P_e\mu_e - P_f\mu_f)\right)\}^2] \]
When $P_e-P_f$ or equivalently $\Sigma_e - \Sigma_f$ has full rank, the set

\[\{x \in \R^d: u^\top(P_e-P_f)x = u^\top(P_e\mu_e - P_f\mu_f)\}\]

is a hyperplane and therefore has probability zero under a nondegenerate Gaussian mixture. Therefore,

\[u^\top M u > 0\],

and $M$ would be invertible.


\end{remark}

\begin{assumption}[Conditions for the KRR minimax rate]\label{assump:krr}

In each of the environments $g \in \{e,f\}$, we assume the following:

\begin{enumerate}[label=(\alph*),leftmargin=2.2em]
  \item \textbf{Conditional Bernstein moments.} There exists $L,B_U>0$ such that, for every integer $m\ge2$,
  \begin{equation}
    \E[|U_i|^m\mid X_i=x,E_i=g]
    \le \frac12m!B_U^2L^{m-2}
    \label{eq:bernstein_moments}
  \end{equation}
 This condition is implied by \ref{assump:model}.

  \item \textbf{Spectral behavior.} We distinguish the following two cases.
  \begin{enumerate}[label=(\roman*),leftmargin=2.2em]
    \item \textbf{Infinite-rank case.} If $f_0$ belongs to an infinite dimensional RKHS, we assume the two-sided eigenvalue decay condition in each environment $g \in \{e,f\}$. In each environment $g$, for some $0<p_g<1$ and constants $0<c_g\le C_g<\infty$,
    \begin{equation}
      c_g j^{-1/p}\le\mu_{g,j}\le C_gj^{-1/p},\qquad j\ge1.
      \label{eq:eigen_decay}
    \end{equation}

    \item \textbf{Finite-rank case.} If $f_0$ belongs to an infinite dimensional RKHS,  for each environment $g \in \{e,f\}$, the lowest eigenvalue of the environment covariance operator is bounded away from zero 
    \begin{equation}
      \mu_{r_g}>0.
      \label{eq:finite_spectral_gap}
    \end{equation}

  \end{enumerate}

  \item \textbf{Source condition.}
  \begin{enumerate}[label=(\roman*),leftmargin=2.2em]
    \item  For some $1<\beta_g\le2$ and $B_g>0$,
    \begin{equation}
      [f_0]_g\in[\F]^{\beta_g},\qquad
      \norm{[f_0]_g}_{[\F]^{\beta_g}_g}\le B_g.
      \label{eq:source_condition}
    \end{equation}
    Equivalently, there exists $[h_0]_g$ with $\norm{[h_0]_g}_{L_2(\nu)}\le B_g$ such that $[f_0]_g=T_g^{{\beta_g}/2}[h_0]_g$.
  \end{enumerate}
\end{enumerate}
\end{assumption}

\begin{remark}
These conditions are similar to other papers in the literature examining optimal rates for KRR estimators \cite{blanchard2018optimal,fischer2020sobolev,lin2020optimal}. In the infinite-rank case, the source condition is necessary. Without one, restricted to estimation of $f_0$ in one environment alone, it can be shown that in an infinite-dimensional RKHS ball alone, there is no uniform RKHS-norm rate. The range $\beta>1$ yields consistency in the first power norm, and hence in $\F$ under \Cref{assump:model}. 
\end{remark}

\begin{remark}
For the finite-rank case, the source condition is automatically satisfied since there are only finitely many positive eigenvalues.
\end{remark}

\begin{assumption}[Local leverage regular variation at $x_0$]\label{assump:local_leverage}
 If $f_0$ belongs to an infinite dimensional RKHS, there exist $A_0>0$ and $0<a<1$ such that
  \begin{equation}
    \Ncal_{x_0}(\eta)\sim A_0\eta^{-a}\qquad(\eta\downarrow0).
    \label{eq:local_regular_variation}
  \end{equation}
  As a result, for every fixed $c>0$,
  \begin{equation}
    \frac{\Ncal_{x_0}(c\eta)}{\Ncal_{x_0}(\eta)}\to c^{-a}.
    \label{eq:local_ratio}
  \end{equation}

\end{assumption}


\begin{assumption}[Query point non-degeneracy]\label{assump:nondegeneracy}
We assume $k_0=(K(x_0,X_1),\ldots,K(x_0,X_n))^\top$ is non-zero almost surely. This implies that the objective function in \Cref{prop:gm_dual} is strictly concave, which is necessary to show $\mu_n^\star=\eta_n^{-1}$ in \Cref{prop:regularized_gamma}.
\end{assumption}


% ============================================================


\section{Methodology}
% ============================================================

Given data from \Cref{assump:model}, the goal is pointwise inference for $\theta_0=f_0(x_0)$. We first estimate the confounding quantities using $I_1$. The procedure for constructing a GM debiasing estimator on $I_2$ then consists of three parts: a regularized pilot estimate, minimum-variance balancing weights, and a one-step residual correction.

For each environment, estimate $\mu_g$ and $\Omega_g$ from the observations in $I_1$, and set $\widehat P_g=\widehat\Omega_g^{-1}$. Write $\varphi_g$ and $\widehat\varphi_g$ for the population and fitted Gaussian densities. Define
\[
 q_g(x)=\frac{\pi_g\varphi_g(x)}{\sum_h\pi_h\varphi_h(x)},\qquad
 \widehat q_g(x)=\frac{\widehat\pi_g\widehat\varphi_g(x)}{\sum_h\widehat\pi_h\widehat\varphi_h(x)}.
\]
For fixed environment proportions, take $\widehat\pi_g=\pi_g$; for i.i.d. mixture sampling, use the first-fold proportions. For $i\in I_1$, let
\[
 \widehat V_i=\widehat P_{E_i}(X_i-\widehat\mu_{E_i})
 -\sum_g\widehat q_g(X_i)\widehat P_g(X_i-\widehat\mu_g).
\]
Then set
\begin{equation}
 \widehat M=\frac1m\sum_{i\in I_1}\widehat V_i\widehat V_i^\top,
 \qquad \widehat h=\frac1m\sum_{i\in I_1}\widehat V_iY_i,
 \qquad \widehat c=(\widehat M+\tau_m I_d)^{-1}\widehat h,
 \label{eq:nuisance_estimation}
\end{equation}
where $\tau_m=0$ or, for example, $\tau_m=m^{-1}$. More generally, $\tau_m=O(m^{-1/2})$ suffices. Estimates are defined arbitrarily on the vanishing-probability events where an inverse does not exist. Define
\[
 \widehat m_g(x):=\widehat c^\top\widehat P_g(x-\widehat\mu_g),
 \qquad \widetilde Y_i:=Y_i-\widehat m_{E_i}(X_i).
\]
The two error variances are estimated from first-fold corrected KRR residuals as described below. All these quantities are held fixed when constructing the GM estimator on $I_2$.

\begin{definition}[GM debiasing estimator]\label{def:gm_estimator}
Fix a query point $x_0\in\X$.

\medskip
\noindent\textbf{Step 1 (KRR pilot).} Compute
\begin{equation}
  \widehat f_{\lambda}
  =\argmin_{f\in\F}
  \left\{
    \frac1n\sum_{i=1}^n(\widetilde Y_i-f(X_i))^2+\lambda\norm{f}_{\F}^2
  \right\}.
  \label{eq:krr_pilot}
\end{equation}
Equivalently, $\widehat\alpha=(K+n\lambda I)^{-1}\widetilde{\bm Y}$ and $\widehat f_\lambda(x)=k_x^\top\widehat\alpha$, where $k_x=(K(x,X_1),\ldots,K(x,X_n))^\top$.

\medskip
\noindent\textbf{Step 2 (GM debiasing weights).} For a tolerance $\gamma\ge0$, solve
\begin{equation}
  \widehat\omega
  \in\argmin_{\omega\in\R^n}\frac1n\omega^\top W\omega
  \quad\text{subject to}\quad
  \norm{K_{x_0}-\frac1n\sum_{i=1}^n\omega_iK_{X_i}}_{\F}\le\gamma.
  \label{eq:gm_primal}
\end{equation}

\medskip
\noindent\textbf{Step 3 (one-step correction).} Define
\begin{equation}
  \widehat\theta_{\mathrm{GM}}
  :=\widehat f_\lambda(x_0)
  +\frac1n\widehat\omega^\top\{\widetilde{\bm Y}-\widehat f_\lambda(\bm X)\}.
  \label{eq:gm_estimator}
\end{equation}
\end{definition}

We have the following identity:
\begin{equation}
  \widehat\theta_{\mathrm{GM}}-\theta_0
  =\frac1n\widehat\omega^\top\bm U
  +\inner{\widehat f_\lambda-f_0}{\Delta_{\widehat\omega}}_{\F}
  +R_{\mathrm{nuis},n},
  \label{eq:gm_decomposition}
\end{equation}
where
\[
 R_{\mathrm{nuis},n}:=\frac1n\sum_{i=1}^n\widehat\omega_i
 \{m_{E_i}(X_i)-\widehat m_{E_i}(X_i)\}.
\]
And thus,
\[
  \abs{\inner{\widehat f_\lambda-f_0}{\Delta_{\widehat\omega}}_{\F}}
  \le\norm{\widehat f_\lambda-f_0}_{\F}\,\norm{\Delta_{\widehat\omega}}_{\F}
  \le \gamma\norm{\widehat f_\lambda-f_0}_{\F}.
\]
The GM primal program in \eqref{eq:gm_primal} therefore directly trades estimated conditional variance, proportional to $\omega^\top W\omega$, against the worst-case RKHS alignment error.

\begin{remark}[When exact alignment is possible]
The condition $\gamma=0$ is equivalent to
\[
  K_{x_0}\in\operatorname{span}\{K_{X_1},\ldots,K_{X_n}\},
\]

In a finite-dimensional RKHS, exact alignment can occur once the sampled feature vectors span the feature space. More generally, if $C_\nu$ has finite rank and $K_{x_0}\in\F$, exact alignment can occur once the sampled kernel sections span $\F$. For the inhomogeneous degree-$q$ polynomial kernel on $\R^p$, the feature-space dimension is $\binom{p+q}{q}$, so exact alignment can occur once $n\geq\binom{p+q}{q}$ and the sampled kernel sections span the feature space (for example, when $\binom{p+q}{q}$ of them are linearly independent). In contrast, for a strictly positive-definite infinite-dimensional kernel under an infinite-rank design and distinct points $x_0,X_1,\ldots,X_n$, the kernel representations $K_{x_0},K_{X_1},\dots,K_{X_n}$ are linearly independent, so exact alignment is impossible with finitely many observations.
\end{remark}

\subsection{Dual formulation and numerical solution}


\begin{proposition}[Dual of the GM debiasing program]
\pdfbookmark[2]{Proposition \theproposition}{proposition.bookmark.\theproposition}
\label{prop:gm_dual}
Let $K$, $k_0$, and $k_{00}$ be as defined above. Suppose Slater's condition holds, or equivalently, there exists a weight vector with strictly smaller alignment error than $\gamma$. In particular, using \eqref{eq:gamma_min_closed}, strict feasibility holds if and only if
\[
  \gamma^2>\gamma_{\min,n}^2
  =k_{00}-k_0^\top K^\dagger k_0.
\]
Then the following statements hold:
\begin{enumerate}[label=(\roman*)]

  \item The primal program \eqref{eq:gm_primal} is equivalent to the one-dimensional concave dual
  \begin{equation}
    \max_{\mu\ge0}
    \left\{
    \mu(k_{00}-\gamma^2)
    -\frac{\mu^2}{n^2}k_0^\top
    \left(\frac1nW+\frac{\mu}{n^2}K\right)^{-1}k_0
    \right\},
    \label{eq:gm_dual}
  \end{equation}
  and the primal weights are recovered by
  \begin{equation}
    \widehat\omega
    =\frac{\mu^\star}{n}
    \left(\frac1nW+\frac{\mu^\star}{n^2}K\right)^{-1}k_0.
    \label{eq:dual_weight_recovery}
  \end{equation}

  \item If $\widetilde K:=W^{-1/2}KW^{-1/2}=U\Lambda U^\top$ with
  $\Lambda=\operatorname{diag}(\lambda_1,\ldots,\lambda_n)$ and
  $a:=U^\top W^{-1/2}k_0$, then the dual objective is
  \begin{equation}
    f(\mu)
    =(k_{00}-\gamma^2)\mu
    -\sum_{j=1}^n\frac{a_j^2\mu^2}{n+\lambda_j\mu},
    \label{eq:dual_eigenbasis}
  \end{equation}
  with
  \begin{align}
    f'(\mu)
    &=k_{00}-\gamma^2-
    \sum_{j=1}^n a_j^2
    \frac{\mu(2n+\lambda_j\mu)}
    {(n+\lambda_j\mu)^2},
    \label{eq:dual_first_derivative}\\
    f''(\mu)
    &=-\sum_{j=1}^n
    a_j^2\frac{2n^2}
    {(n+\lambda_j\mu)^3}
    \leq 0.
    \label{eq:dual_second_derivative}
  \end{align}

\end{enumerate}
\end{proposition}



\hyperref[proof:gm_dual]{A proof of Proposition~\ref*{prop:gm_dual}} can be found in \Cref{sec:additional_proofs}. 

To solve the GM dual program as stated in \eqref{eq:gm_dual}, a Newton--Raphson algorithm can be implemented as outlined in Algorithm~\ref{alg:newton_dual}.

\begin{algorithm}[ht]
\caption{Newton method for the GM dual}
\label{alg:newton_dual}
\begin{algorithmic}[1]
\Require Eigendecomposition $W^{-1/2}KW^{-1/2}=U\Lambda U^\top$, $a=U^\top W^{-1/2}k_0$, $k_{00}$, $\gamma^2$, tolerance $\mathrm{tol}$.
\If{$k_{00}-\gamma^2\le0$}
  \State \Return $\widehat\mu=0$.
\EndIf
\State Set $\mu_L\gets0$ and obtain $\mu_U$ with $f'(\mu_U)<0$ by repeated doubling.
\State Set $\mu\gets\mu_U/2$.
\Repeat
  \State Update the bracket using the sign of $f'(\mu)$.
  \State Set $\mu_{\mathrm{new}}\gets\mu-f'(\mu)/f''(\mu)$.
  \If{$\mu_{\mathrm{new}}\notin(\mu_L,\mu_U)$ or is non-finite}
    \State $\mu_{\mathrm{new}}\gets(\mu_L+\mu_U)/2$.
  \EndIf
  \State $\mu\gets\mu_{\mathrm{new}}$.
\Until{$|f'(\mu)|<\mathrm{tol}$}
\State Recover $\widehat\omega=\mu n\,W^{-1/2}U\{a_j/(n+\mu\lambda_j)\}_{j=1}^n$.
\end{algorithmic}
\end{algorithm}

\subsection{Variance estimation and studentization}
Under conditional homoskedasticity within each environment, the leading noise variance is
\begin{equation}
 \operatorname{Var}\!\left(\frac1n\widehat\omega^\top\bm U
 \,\middle|\,I_1,\bm X,\bm E\right)
 =\frac1{n^2}\sum_{g\in\{e,f\}}\sigma_{U,g}^2
      \sum_{i:E_i=g}\widehat\omega_i^2.
 \label{eq:conditional_variance}
\end{equation}
In the implementation, the two variances are estimated from corrected KRR residuals using only $I_1$. Fit $\widehat f_{\lambda_m}^{(1)}$ on the first-fold corrected responses, with $\lambda_m\asymp m^{-1/(\beta+p)}$ in the infinite-rank case and $\lambda_m=O(m^{-1/2})$ in the finite-rank case. An estimator for $\sigma_{U,g}^2$ is thus
\begin{equation}
 \widehat\sigma_{U,g}^2
 :=\max\left\{m^{-1},\frac1{m_g}\sum_{i\in I_1:E_i=g}
 \{Y_i-\widehat m_g(X_i)-\widehat f_{\lambda_m}^{(1)}(X_i)\}^2\right\}.
 \label{eq:sigma_hat}
\end{equation}
The lower truncation ensures that $W$ is positive definite. It has no asymptotic effect, since $\sigma_{U,g}^2>0$.

The asymptotic variance of our GM estimator can then be estimated by
\begin{equation}
 \widehat{\operatorname{Var}}_{\mathrm{asy}}(\widehat\theta_{\mathrm{GM}})
 :=\frac1{n^2}\sum_{g\in\{e,f\}}\widehat\sigma_{U,g}^2
     \sum_{i:E_i=g}\widehat\omega_i^2
 =\frac1n\widehat V_n,
 \label{eq:implemented_variance}
\end{equation}
where $\widehat V_n:=n^{-1}\widehat\omega^\top W\widehat\omega$. The standard error is $\sqrt{\widehat V_n/n}$. The same observations in $I_1$ may be used for the nuisance estimates, the first-fold KRR fit, and the residual variance estimates; consistency is proved in \Cref{prop:noise_variances}.

\clearpage
\newpage
\section{Main Proofs}

\begin{theorem}[CLT for the RKHS Gauss--Markov estimator]\label{thm:clt}
\pdfbookmark[2]{Theorem \thetheorem}{theorem.bookmark.\thetheorem}
Suppose Assumptions~\ref{assump:model}--\ref{assump:nondegeneracy} hold, $r_\nu=\infty$, and set the tolerance parameter
\[
 \gamma_n=\norm{\eta_n(D_{\nu,n}+\eta_n I)^{-1}K_{x_0}}_{\F},
 \qquad \eta_n\asymp(\log n)^2/n.
\]
Choose the KRR pilot regularization as $\lambda_n\asymp n^{-1/(\beta+p)}$, and construct the nuisance and variance estimates from $I_1$ as in \Cref{def:gm_estimator} and \eqref{eq:nuisance_estimation}--\eqref{eq:sigma_hat}, with $\tau_m=O(m^{-1/2})$. Then
\begin{equation}
 \frac{\sqrt n\,(\widehat\theta_{\mathrm{GM}}-\theta_0)}{\sqrt{\widehat V_n}}
 \overset d\longrightarrow N(0,1),
 \qquad \widehat V_n:=\frac1n\sum_{i=1}^n\widehat\sigma_{U,E_i}^2\widehat\omega_i^2.
 \label{eq:clt}
\end{equation}
\end{theorem}

\begin{proof}
By \eqref{eq:gm_decomposition},
\begin{align}
 \frac{\sqrt n(\widehat\theta_{\mathrm{GM}}-\theta_0)}{\sqrt{\widehat V_n}}
 &=\underbrace{\frac{\widehat\omega^\top\bm U}{\sqrt{n\widehat V_n}}}_{A_n}
 +\underbrace{\frac{\sqrt n}{\sqrt{\widehat V_n}}
   \inner{\widehat f_{\lambda_n}-f_0}{\Delta_{\widehat\omega}}_{\F}}_{B_n}
 +\underbrace{\frac{\sqrt n\,R_{\mathrm{nuis},n}}{\sqrt{\widehat V_n}}}_{C_n}.
 \label{eq:clt_decomposition}
\end{align}
We first show that the bias term $B_n$ vanishes as $n\to\infty$. By feasibility and Cauchy--Schwarz,
\[
 \abs{\inner{\widehat f_{\lambda_n}-f_0}{\Delta_{\widehat\omega}}_{\F}}
 \leq\gamma_n\norm{\widehat f_{\lambda_n}-f_0}_{\F}.
\]
\Cref{prop:regularized_gamma} gives $\gamma_n^2=O_p(\eta_n^{1-a})$.
\Cref{prop:variance_consistency} gives $\widehat V_n/V_n\to_p1$ and $V_n\asymp\eta_n^{-a}$. Therefore, since $n\eta_n\asymp(\log n)^2$,
\begin{align*}
 \frac{\sqrt n}{\sqrt{\widehat V_n}}
 \abs{\inner{\widehat f_{\lambda_n}-f_0}{\Delta_{\widehat\omega}}_{\F}}
 &\leq\left(\frac{V_n}{\widehat V_n}\right)^{1/2}
 \left(\frac{n\gamma_n^2}{V_n}\right)^{1/2}\norm{\widehat f_{\lambda_n}-f_0}_{\F}\\
 &=O_p(\sqrt{n\eta_n})\norm{\widehat f_{\lambda_n}-f_0}_{\F}\\
 &=O_p\!\left((\log n)n^{-(\beta-1)/\{2(\beta+p)\}}\right)=o_p(1),
\end{align*}
where the last step is \Cref{thm:krr_rate} and $\beta>1$.

Next, we show that the nuisance term $C_n$ vanishes. Write
\[
 a_g:=P_gc,\qquad b_g:=-c^\top P_g\mu_g,
 \qquad m_g(x)=a_g^\top x+b_g,
\]
and use hats for the corresponding plug-in estimates. By \Cref{prop:nuisance_rates}, their errors are $O_p(m^{-1/2})$. We have the exact identity
\begin{align*}
 R_{\mathrm{nuis},n}
 &=\sum_g(a_g-\widehat a_g)^\top
 \frac1n\sum_{i=1}^n\widehat\omega_i\mathbf1\{E_i=g\}X_i\\
 &\quad+\sum_g(b_g-\widehat b_g)
 \frac1n\sum_{i=1}^n\widehat\omega_i\mathbf1\{E_i=g\}.
\end{align*}
Each of the finitely many features $\mathbf1\{E=g\}X_j$ and $\mathbf1\{E=g\}$ is square-integrable. \Cref{lem:weighted_moments} therefore implies
\[
 R_{\mathrm{nuis},n}=o_p\!\left(m^{-1/2}\sqrt{V_n}\right),\qquad
 C_n=o_p\!\left(\sqrt{n/m}\right)=o_p(1).
\]
This proves the required nuisance remainder bound from the generating model and first-fold estimation; it is not an additional assumption.

Next, we show that, conditional on the first fold and the second-fold covariates and environment labels, $A_n$ converges in distribution to $N(0,1)$ in probability. Denote this information by $\mathcal G_n$. The variables $U_i$ are conditionally independent and centered. Their conditional variance gives
\[
 V_n^{\mathrm{noise}}:=\frac1n\sum_{i=1}^n\sigma_{U,E_i}^2\widehat\omega_i^2,
 \qquad
 \abs{\frac{V_n^{\mathrm{noise}}}{\widehat V_n}-1}
 \leq\max_g\abs{\frac{\sigma_{U,g}^2}{\widehat\sigma_{U,g}^2}-1}
 \overset p\longrightarrow0
\]
by \Cref{prop:noise_variances}. The conditional Bernstein condition gives a uniform fourth-moment bound. Thus
\[
 \frac{\sum_i\E[(\widehat\omega_iU_i)^4\mid\mathcal G_n]}
 {\{\sum_i\widehat\omega_i^2\sigma_{U,E_i}^2\}^2}
 \leq C\frac{\max_i\widehat\omega_i^2}{\sum_i\widehat\omega_i^2}
 \overset p\longrightarrow0
\]
by \Cref{prop:max_sum}. Hence the conditional Lyapunov condition is fulfilled, and the conditional CLT gives convergence of the noise sum normalized by $\sqrt{nV_n^{\mathrm{noise}}}$ to $N(0,1)$ in probability. Ratio consistency replaces $V_n^{\mathrm{noise}}$ by $\widehat V_n$.

Finally, we employ an argument using bounded continuous functions to upgrade our result to unconditional convergence in distribution. For every bounded continuous function $g$,
\[
 \E[g(A_n)\mid\mathcal G_n]\overset p\longrightarrow\E[g(Z)],\qquad Z\sim N(0,1).
\]
The conditional expectations are uniformly bounded, so convergence in probability to the constant also implies convergence in $L_1$. Therefore $\E[g(A_n)]\to\E[g(Z)]$. Combining with $B_n=C_n=o_p(1)$ and Slutsky's theorem proves the result.
\end{proof}

\begin{theorem}[Minimax RKHS-norm rate for KRR]\label{thm:krr_rate}
\pdfbookmark[2]{Theorem \thetheorem}{theorem.bookmark.\thetheorem}
Suppose Assumptions~\ref{assump:model} and \ref{assump:krr} hold, and use the first-fold nuisance estimates with $\tau_m=O(m^{-1/2})$. Then the following statements hold.

\begin{enumerate}[label=(\roman*),leftmargin=2.2em]
  \item \textbf{Infinite-rank case.} Suppose $r_\nu=\infty$, $1<\beta\le2$, and let
  \[
    \lambda_n\asymp n^{-1/(\beta+p)}.
  \]
  Then
  \begin{equation}
    \norm{\widehat f_{\lambda_n}-f_0}_{\F}
    =O_p\!\left(n^{-(\beta-1)/\{2(\beta+p)\}}\right).
    \label{eq:krr_rate}
  \end{equation}
  The polynomial exponent is the same as in the oracle regression problem satisfying the two-sided eigenvalue and source conditions.

  \item \textbf{Finite-rank case.} Suppose $r_\nu<\infty$ and \eqref{eq:finite_spectral_gap} holds. If $\lambda_n\to0$, then
  \begin{equation}
    \norm{\widehat f_{\lambda_n}-f_0}_{\F}
    =O_p\!\left(n^{-1/2}+\lambda_n\right).
    \label{eq:krr_rate_finite}
  \end{equation}
  In particular, if $\lambda_n=O(n^{-1/2})$, then
  \begin{equation}
    \norm{\widehat f_{\lambda_n}-f_0}_{\F}
    =O_p(n^{-1/2}).
    \label{eq:krr_rate_finite_parametric}
  \end{equation}
\end{enumerate}
\end{theorem}

\begin{proof}
We treat the two cases separately.

\medskip
\noindent\textbf{Infinite-rank case.}
Let $\widehat f_{\lambda_n}^{\circ}$ be the KRR fit using oracle responses $f_0(X_i)+U_i$. \Cref{lem:first_power_isometry} identifies the RKHS norm on $\F$ with the first power-space norm. The usual oracle rate agrees with Theorems 1 and 2 of \citet{fischer2020sobolev}. We give a direct bound which also applies to the two fixed strata.

Put $s=(\beta-1)/2$, so the source condition gives $f_0=C_\nu^sv$ for a fixed $v\in\F$. By \Cref{lem:operator_bernstein}, $C_{\nu,n}+\lambda_n I\succeq(1-o_p(1))(C_\nu+\lambda_n I)$ since $n\lambda_n/\log n\to\infty$. The KRR normal equation and the spectral theorem yield
\begin{align*}
 \norm{\widehat f_{\lambda_n}^{\circ}-f_0}_{\F}
 &\leq\norm{\lambda_n(C_{\nu,n}+\lambda_n I)^{-1}C_\nu^sv}_{\F}\\
 &\quad+\norm{(C_{\nu,n}+\lambda_n I)^{-1}\frac1n\sum_iU_iK_{X_i}}_{\F}\\
 &=O_p\!\left(\lambda_n^s+(n\lambda_n^{1+p})^{-1/2}\right).
\end{align*}
For the bias term, use $\norm{\lambda_n(C_{\nu,n}+\lambda_n I)^{-1}C_\nu^sv}\lesssim\sqrt{\lambda_n}\norm{(C_\nu+\lambda_n I)^{-1/2}C_\nu^sv}\lesssim\lambda_n^s\norm v$. For the noise term, regularized operator comparison and
\[
 \E\left\|(C_\nu+\lambda_n I)^{-1/2}\frac1n\sum_iU_iK_{X_i}\right\|_{\F}^2
 \leq\frac{\max_g\sigma_{U,g}^2}{n}\Ncal_\nu(\lambda_n)
 \lesssim\frac{\lambda_n^{-p}}n
\]
give the displayed bound. The chosen $\lambda_n$ gives the stated oracle rate.

Let $r_i=m_{E_i}(X_i)-\widehat m_{E_i}(X_i)$. \Cref{prop:nuisance_rates} and the Gaussian second moments give $n^{-1}\sum_i r_i^2=O_p(m^{-1})$. The deterministic KRR perturbation bound is
\[
 \norm{\widehat f_{\lambda_n}-\widehat f_{\lambda_n}^{\circ}}_{\F}
 \leq\frac1{2\sqrt{\lambda_n}}\left(\frac1n\sum_i r_i^2\right)^{1/2}
 =O_p((m\lambda_n)^{-1/2}),
\]
since $\sup_{t\geq0}\sqrt t/(t+\lambda_n)=1/(2\sqrt{\lambda_n})$. This rate is faster than the oracle rate when $m\asymp n$, proving \eqref{eq:krr_rate}. The same proof applies to the first-fold KRR fit using its own nuisance estimates: the affine coefficient errors factor out of empirical second moments, so no independence within $I_1$ is needed for this perturbation bound.

\medskip
\noindent\textbf{Finite-rank case.}
The KRR normal equation gives
\[
  (C_{\nu,n}+\lambda_n I)\widehat f_{\lambda_n}
  =
  C_{\nu,n}f_0+\frac1n\sum_{i=1}^n(U_i+r_i)K_{X_i}.
\]

Subtracting $(C_{\nu,n}+\lambda_n I) f_0$  on both sides gives

\begin{equation}
  \widehat f_{\lambda_n}-f_0
  =
  (C_{\nu,n}+\lambda_n I)^{-1}
  \left(\frac1n\sum_{i=1}^n(U_i+r_i)K_{X_i}\right)
  -
  \lambda_n(C_{\nu,n}+\lambda_n I)^{-1}f_0,
  \label{eq:krr_finite_decomp}
\end{equation}
Since $r_\nu<\infty$ and $C_\nu\succeqL\mu_{\min}^+I$, boundedness of the kernel gives
\[
  \E\norm{C_{\nu,n}-C_\nu}_{\HS}^2
  \leq
  \frac1n\E\norm{K_X\otimes K_X-C_\nu}_{\HS}^2
   \leq
  \frac1n\E\norm{K_X\otimes K_X}_{\HS}^2
  =O(n^{-1}).
\]
Since $\norm{\cdot}_{\op}\le\norm{\cdot}_{\HS}$, using Chebyshev's inequality, we get
\[
  \norm{C_{\nu,n}-C_\nu}_{\op}=O_p(n^{-1/2}),
\]

Note that 
\[C_{\nu,n} \succeq C_\nu - \|C_{\nu,n} - C_{\nu}\|_{\op} I \succeq \mu^+_{\min}I - \|C_{\nu,n} - C_{\nu}\|_{\op} I\] 

Therefore, as $n \to \infty$, we have 
\[
  \Pp\!\left(
  C_{\nu,n}\succeqL\frac{\mu_{\min}^+}{2}I
  \right)\longrightarrow1.
\]
On this event,
\[
  \norm{(C_{\nu,n}+\lambda_n I)^{-1}}_{\op}
  \le\frac{2}{\mu_{\min}^+}.
\]
Moreover, using $\E[U_i\mid X_i,E_i]=0$, $\E[U_i^2\mid X_i,E_i]\leq\max_g\sigma_{U,g}^2$, and $\norm{K_{X_i}}_{\F}\le\kappa$,
\[
  \E\left\|
  \frac1n\sum_{i=1}^nU_iK_{X_i}
  \right\|_{\F}^2
  \le
  \frac{(\max_g\sigma_{U,g}^2)\kappa^2}{n}.
\]
Also,
\[
 \norm{\frac1n\sum_i r_iK_{X_i}}_{\F}
 \leq\kappa\left(\frac1n\sum_i r_i^2\right)^{1/2}=O_p(m^{-1/2}).
\]
Thus the first term in \eqref{eq:krr_finite_decomp} is $O_p(n^{-1/2})$ and the second is $O_p(\lambda_n)$, proving \eqref{eq:krr_rate_finite}. Letting $\lambda_n=O(n^{-1/2})$, we thus prove our claim.
\end{proof}


\begin{proposition}
[Rate from regularized choice of $\gamma_n$]
\pdfbookmark[2]{Proposition \theproposition}{proposition.bookmark.\theproposition}
\label{prop:regularized_gamma}
Suppose Assumptions~\ref{assump:model}--\ref{assump:nondegeneracy} hold. For $\eta>0$, define
\[
  R_n^2(\eta)
  :=\norm{\eta(D_{\nu,n}+\eta I)^{-1}K_{x_0}}_{\F}^2.
\]
Choose $\gamma_n^2=R_n^2(\eta_n)$. Then the following statements hold:
\begin{enumerate}[label=(\roman*)]
  \item Slater's condition holds, and the unique dual optimum is
  \[
    \mu_n^\star=\eta_n^{-1}.
  \]

  \item If $\eta_n\asymp (\log n)^2/n$ and $r_\nu=\infty$, then
  \[
    \gamma_n=O_p\!\left(n^{-(1-a)/2}(\log n)^{1-a}\right),
    \qquad
    \mu_n^\star\asymp \frac{n}{(\log n)^2}.
  \]

  \item If $\eta_n\asymp n^{-1}$ and $r_\nu<\infty$, then
  \[
    \gamma_n=O_p(n^{-1}),
    \qquad
    \mu_n^\star\asymp n.
  \]
\end{enumerate}
\end{proposition}




\begin{proof}

First note from \Cref{lem:weight_operator_form} that the weights from the dual solution in \Cref{prop:gm_dual} have the following representation in terms of the empirical covariance operator $D_{\nu,n}$:

\[
  \widehat\omega(\eta_n)=W^{-1}S_n(D_{\nu,n}+\eta_n I)^{-1}K_{x_0}.
\]

The RKHS alignment error from the GM primal program~\eqref{eq:gm_primal} for such weights is 

\begin{align*}
   \norm{K_{x_0} - \frac{1}{n} \sum_{i=1}^n \widehat\omega_i K_{X_i}}_{\F} &= \norm{K_{x_0} - S_n^*W^{-1}S_n(D_{\nu,n} + \eta_n I)^{-1}K_{x_0}}_{\F}\\
   &=\norm{(I - D_{\nu,n}(D_{\nu,n} + \eta_n I)^{-1})K_{x_0}}_{\F}\\
   &= \norm{\eta_n(D_{\nu,n}+\eta_n I)^{-1}K_{x_0}}_{\F} \\
   &=R_n(\eta_n)=:\gamma_n
\end{align*}

This implies that it suffices to prove that $R_n(\eta)$ is strictly increasing since the weights $\widehat\omega(\eta_n/2)$ result in an alignment error less than $\gamma_n$ and thus Slater's condition applies. 

$D_{\nu,n}$ is a compact, self-adjoint, and positive semidefinite operator. By the spectral theorem, it can be decomposed as the sum

\[D_{\nu,n}=\sum_{j=1}^{r_n}\widehat{\mu}_{j,n}\,\widehat e_{j,n}\otimes\widehat e_{j,n}.\]

where $\widehat{\mu}_{j,n}>0$ and $\widehat e_{j,n}$ are the non-zero eigenvalues and corresponding orthonormal eigenvectors of $D_{\nu,n}$. 
We can first decompose $R^2_n(\eta)$ as the sum of two parts, the projection onto $\ran(D_{\nu,n})$ and its projection onto $\ker(D_{\nu,n})$:

\[K_{x_0} =
\sum_{j=1}^{r_n}
\left\langle K_{x_0},\widehat e_{j,n}\right\rangle_{\F}
\widehat e_{j,n}
+
P_{\ker(D_{\nu,n})}K_{x_0}\]

Left-multiplying by $\eta(D_{\nu,n} + \eta I)^{-1}$ gives

\begin{align*}
\eta(D_{\nu,n}+\eta I)^{-1}K_{x_0}
&=
\sum_{j=1}^{r_n}
\frac{\eta}{\widehat\mu_{j,n}+\eta}
\left\langle K_{x_0},\widehat e_{j,n}\right\rangle_{\F}
\widehat e_{j,n}
+
P_{\ker(D_{\nu,n})}K_{x_0},
\\\implies
\left\|
\eta(D_{\nu,n}+\eta I)^{-1}K_{x_0}
\right\|_{\F}^{2}
&=
\sum_{j=1}^{r_n}
\left(
\frac{\eta}{\widehat\mu_{j,n}+\eta}
\right)^2
\left|
\left\langle K_{x_0},\widehat e_{j,n}\right\rangle_{\F}
\right|^2
+
\left\|
P_{\ker(D_{\nu,n})}K_{x_0}
\right\|_{\F}^{2}.
\end{align*}


For any constant $c>0$, the function $x\mapsto x/(x+c)$ is strictly increasing. This implies that the map $\eta\mapsto R_n^2(\eta)$ is strictly increasing when $k_0\neq0$. 

Hence $\widehat\omega(\eta_n/2)$ is strictly feasible for the tolerance $R_n(\eta_n)$, proving Slater.

For the second part, using notation from \Cref{prop:gm_dual}, note first that $R_n^2(\eta_n)$ has this alternate representation:

\begin{align*}
R_n^2(\eta_n)
&:=\left\|\eta_n(D_{\nu,n}+\eta_n I)^{-1}K_{x_0}\right\|_{\F}^2\\
&=k_{00}-\sum_{j=1}^na_j^2
\frac{\lambda_j+2n\eta_n}
{(\lambda_j+n\eta_n)^2}
\end{align*}


From \Cref{prop:gm_dual}, using \eqref{eq:dual_first_derivative}, we get

\begin{align*}
f'\!\left(\frac{1}{\eta_n}\right)
&=
k_{00}-\gamma_n^2
-\sum_{j=1}^na_j^2\frac{\frac{1}{\eta_n}\left(2n+\frac{\lambda_j}{\eta_n}\right)
}{
\left(n+\frac{\lambda_j}{\eta_n}\right)^2}
\\
&=
k_{00}-\gamma_n^2-\sum_{j=1}^na_j^2\frac{\lambda_j+2n\eta_n}
{(\lambda_j+n\eta_n)^2}\\
&=\left[k_{00}-\sum_{j=1}^na_j^2\frac{\lambda_j+2n\eta_n}{(\lambda_j+n\eta_n)^2}\right]-\gamma_n^2
\\
&=R_n^2(\eta_n)-\gamma_n^2 = 0.
\end{align*}



By \Cref{assump:nondegeneracy}, $k_0\neq0$, so \eqref{eq:dual_second_derivative} is strictly negative for all $\mu\ge0$ and the dual objective is strictly concave. Strict concavity implies that there is at most one stationary point. Hence our optimal solution must be  $\mu_n^\star=1/\eta_n$.

Finally, we obtain the rate for our choice of $\gamma_n$. We distinguish the infinite dimensional and finite dimensional cases. 

\noindent\textbf{Infinite-rank case.}

Define
\[
  E_n:=(D_\nu+\eta_nI)^{-1/2}(D_{\nu,n}-D_\nu)(D_\nu+\eta_nI)^{-1/2}.
\]


By \Cref{lem:precision_concentration}, we have the event $\{\norm{E_n}_{\op}\leq c<1\}$ which has probability tending to $1$ as $n\to\infty$. For the remainder of this case, we work on this good event.

First, we have the following expression for $D_{\nu,n} + \eta_n I$:
\[
D_{\nu,n}+\eta_nI
=(D_\nu+\eta_nI)^{1/2}(I+E_n)(D_\nu+\eta_nI)^{1/2}.
\]
Using the good event, we have
\begin{align*}
    \|E_n\|_{\op} \leq c < 1
    &\implies
    -cI \preceqL E_n \preceqL cI
    \\
    &\implies
    (1-c)I \preceqL I+E_n
    \\
    &\implies
    (1-c)(D_\nu+\eta_n I)
    \preceqL
    D_{\nu,n}+\eta_n I.
\end{align*}
Since both operators are strictly positive for $\eta_n>0$, inversion reverses the Loewner order. Hence
\[
    (D_{\nu,n}+\eta_n I)^{-1}
    \preceqL
    \frac{1}{1-c}(D_\nu+\eta_n I)^{-1}.
\]

\medskip
Since $D_{\nu,n}\succeq0$,
\[
  (D_{\nu,n}+\eta_n I)^{-2}
  \preceqL
  \eta_n^{-1}(D_{\nu,n}+\eta_n I)^{-1}.
\]
Therefore
\begin{align*}
  \gamma_n^2
  &= \norm{\eta_n (D_{\nu,n} + \eta_n I)^{-1} K_{x_0}}_{\F}^2\\
  &=\eta_n^2\inner{K_{x_0}}{(D_{\nu,n}+\eta_n I)^{-2} K_{x_0}}_{\F}\\
  &\le\eta_n\inner{K_{x_0}}{(D_{\nu,n}+\eta_nI)^{-1}K_{x_0}}_{\F}\\
  &\le\frac{\eta_n}{1-c}\Ncal_{D,x_0}(\eta_n)
  =O(\eta_n^{1-a}).
\end{align*}
Here \Cref{lem:weighted_scale} gives $\Ncal_{D,x_0}(\eta_n)\asymp\eta_n^{-a}$. Since the good event has probability tending to one and $\eta_n\asymp (\log n)^2/n$, this gives
\[
  \gamma_n=O_p\!\left(n^{-(1-a)/2}(\log n)^{1-a}\right).
\]

\medskip
\noindent\textbf{Finite-rank case.}
Suppose $r_\nu<\infty$. Boundedness of the kernel and consistency of the first-fold precisions imply $\norm{D_{\nu,n}-D_\nu}_{\op}=o_p(1)$.

Therefore, with probability tending to one, we have
\[
  D_{\nu,n}\succeqL\frac{\rho_{\min}^+}{2}I.
\]
On this event $\{ D_{\nu,n}\succeqL\frac{\rho_{\min}^+}{2}I\}$,
\begin{align*}
  \gamma_n
  &=
  \eta_n\norm{(D_{\nu,n}+\eta_n I)^{-1}K_{x_0}}_{\F}\\
  &\le
  \frac{2\eta_n}{\rho_{\min}^+}\norm{K_{x_0}}_{\F}
  =O(\eta_n).
\end{align*}
Thus, when $\eta_n\asymp n^{-1}$,
\[
  \gamma_n=O_p(n^{-1}).
\]

In both cases, $\mu_n^\star=\eta_n^{-1}$ and therefore $\mu_n^\star\asymp n/(\log n)^2$ when $\eta_n\asymp (\log n)^2/n$, and $\mu_n^\star\asymp n$ when $\eta_n\asymp n^{-1}$.

\end{proof}



\begin{proposition}[Ratio consistency of the variance]\label{prop:variance_consistency}
\pdfbookmark[2]{Proposition \theproposition}{proposition.bookmark.\theproposition}
Suppose Assumptions~\ref{assump:model}--\ref{assump:local_leverage} hold. Let $\eta_n\asymp (\log n)^2/n$ if $r_\nu=\infty$ and $\eta_n\asymp n^{-1}$ if $r_\nu<\infty$, and let $\widehat\omega$ be the weights from \Cref{prop:regularized_gamma}. Define
\begin{align}
  \widehat V_n&:=\frac1n\widehat\omega^\top W\widehat\omega,
  \label{eq:Vhat_def}\\
  V_n&:=\inner{(D_\nu+\eta_nI)^{-1}K_{x_0}}
  {D_\nu(D_\nu+\eta_nI)^{-1}K_{x_0}}_{\F}.
  \label{eq:V_def}
\end{align}
Then
\begin{equation}
  \frac{\widehat V_n}{V_n}\overset p\longrightarrow1.
  \label{eq:variance_consistency}
\end{equation}
Moreover, the variance scale depends on the spectral regime:
\begin{enumerate}[label=(\roman*),leftmargin=2.2em]
  \item If $r_\nu=\infty$, then
  \begin{equation}
    V_n\asymp\Ncal_{D,x_0}(\eta_n)\asymp\Ncal_{x_0}(\eta_n)\asymp\eta_n^{-a}\asymp \frac{n^a}{(\log n)^{2a}}.
    \label{eq:V_scale}
  \end{equation}

  \item If $r_\nu<\infty$, then
  \begin{equation}
    V_n
    \longrightarrow
    V_0
    :=
    \inner{K_{x_0}}{D_\nu^\dagger K_{x_0}}_{\F}
    \in(0,\infty),
    \label{eq:V_scale_finite}
  \end{equation}
  and
  \[
    \Ncal_{D,x_0}(\eta_n)\longrightarrow V_0.
  \]
\end{enumerate}
\end{proposition}


\begin{proof}

As in \Cref{lem:precision_concentration}, set
\begin{equation}
  E_n:=(D_\nu+\eta_n I)^{-1/2}(D_{\nu,n}-D_\nu)(D_\nu+\eta_n I)^{-1/2}.
  \label{eq:variance_consistency_E_n}
\end{equation}
\Cref{lem:precision_concentration} provides the event
$\{\norm{E_n}_{\op}\leq c<1\}$, which has probability tending to one as $n\to\infty$.
For the remainder of this proof, we work on this good event.

First, we reexpress $\widehat V_n$ in terms of the empirical precision-weighted covariance operator $D_{\nu,n}$. \Cref{lem:weight_operator_form} yields
\[
 \widehat\omega=W^{-1}S_n(D_{\nu,n}+\eta_n I)^{-1}K_{x_0}.
\]
From \Cref{defn:sampling_operator}, we get
\begin{align*}
 \widehat V_n
 &=\frac1n\widehat\omega^\top W\widehat\omega\\
 &=\inner{(D_{\nu,n}+\eta_n I)^{-1}K_{x_0}}
 {D_{\nu,n}(D_{\nu,n}+\eta_n I)^{-1}K_{x_0}}_{\F}.
\end{align*}

Set
\[
  h_n:=(D_\nu+\eta_n I)^{-1/2}K_{x_0},\qquad
  B_n:=(D_\nu+\eta_n I)^{-1/2}D_\nu(D_\nu+\eta_n I)^{-1/2},
\]
and write
\[
  R_n:=(I+E_n)^{-1}.
\]
Rearranging \eqref{eq:variance_consistency_E_n} gives
\begin{align*}
  D_{\nu,n}+\eta_n I
  &=(D_\nu+\eta_n I)^{1/2}(I+E_n)(D_\nu+\eta_n I)^{1/2},\\
  D_{\nu,n}
  &=(D_\nu+\eta_n I)^{1/2}(B_n+E_n)(D_\nu+\eta_n I)^{1/2}.
\end{align*}
Therefore,
\[
  \widehat V_n
  =\inner{h_n}{R_n(B_n+E_n)R_nh_n}_{\F}.
\]
In contrast, the population variance has the form
\[
  V_n=\inner{h_n}{B_nh_n}_{\F},\qquad
  \norm{h_n}_{\F}^2=\Ncal_{D,x_0}(\eta_n).
\]

On the good event,
\[
  \norm{R_n}_{\op}\leq\frac{1}{1-c},
  \qquad
  R_n-I=-R_nE_n,
  \qquad
  \norm{R_n-I}_{\op}\leq\frac{\norm{E_n}_{\op}}{1-c}.
\]
Moreover,
\[
  0\preceqL B_n=I-\eta_n(D_\nu+\eta_n I)^{-1}\preceqL I,
  \qquad \norm{B_n}_{\op}\leq1.
\]
Using
\[
  R_n(B_n+E_n)R_n-B_n
  =(R_n-I)B_nR_n+B_n(R_n-I)+R_nE_nR_n,
\]
we obtain
\begin{align*}
  \norm{R_n(B_n+E_n)R_n-B_n}_{\op}
  &\leq \norm{R_n-I}_{\op}\norm{R_n}_{\op}
  +\norm{R_n-I}_{\op}
  +\norm{R_n}_{\op}^2\norm{E_n}_{\op}\\
  &\leq C_c\norm{E_n}_{\op},
\end{align*}
where $C_c>0$ depends only on $c$. Hence,
\begin{align*}
  \abs{\widehat V_n-V_n}
  &\leq C_c\Ncal_{D,x_0}(\eta_n)\norm{E_n}_{\op},\\
  \abs{\frac{\widehat V_n}{V_n}-1}
  &\leq C_c\frac{\Ncal_{D,x_0}(\eta_n)}{V_n}\norm{E_n}_{\op}.
\end{align*}
It therefore suffices to prove $\frac{\Ncal_{D,x_0}(\eta_n)}{V_n}=O(1)$. We distinguish the two spectral regimes.

\medskip
\noindent\textbf{Infinite-rank case.}
\Cref{lem:weighted_scale} gives
\[
 V_n\asymp\Ncal_{D,x_0}(\eta_n)\asymp\Ncal_{x_0}(\eta_n)\asymp\eta_n^{-a}.
\]
In particular, $\Ncal_{D,x_0}(\eta_n)/V_n=O(1)$. This step uses the local leverage assumption for $C_\nu$ and comparison with $D_\nu$; no separate regular variation assumption on $D_\nu$ is needed.

\medskip
\noindent\textbf{Finite-rank case.}
Let $b_j:=\inner{K_{x_0}}{\psi_j}_{\F}$. If $r_\nu<\infty$, we immediately observe that as $\eta_n \to 0$
\begin{align*}
V_n
&=
\sum_{j=1}^{r_\nu}
\frac{\rho_j b_j^2}{(\rho_j+\eta_n)^2}
\longrightarrow
\sum_{j=1}^{r_\nu}\frac{b_j^2}{\rho_j},\\
\Ncal_{D,x_0}(\eta_n)
&=
\sum_{j=1}^{r_\nu}\frac{b_j^2}{\rho_j+\eta_n}
\longrightarrow
\sum_{j=1}^{r_\nu}\frac{b_j^2}{\rho_j}.
\end{align*}
The common limit equals
\[
  \inner{K_{x_0}}{D_\nu^\dagger K_{x_0}}_{\F}
  =V_0\in(0,\infty).
\]
Hence $\Ncal_{D,x_0}(\eta_n)/V_n\to1$.

In either case, $\Ncal_{D,x_0}(\eta_n)/V_n=O(1)$. Using $\norm{E_n}_{\op}=o_p(1)$, the ratio of the empirical variance to the population variance thus converges to 1.

\end{proof}

\begin{proposition}[Vanishing max-to-sum ratio]\label{prop:max_sum}
\pdfbookmark[2]{Proposition \theproposition}{proposition.bookmark.\theproposition}
Suppose Assumptions~\ref{assump:model}--\ref{assump:local_leverage} hold. Let $\eta_n\asymp (\log n)^2/n$ if $r_\nu=\infty$ and $\eta_n\asymp n^{-1}$ if $r_\nu<\infty$, and let $\widehat\omega$ be the weights from \Cref{prop:regularized_gamma}. Then
\begin{equation}
  \frac{\max_{1\le i\le n}\widehat\omega_i^2}
  {\sum_{i=1}^n\widehat\omega_i^2}
  \overset p\longrightarrow0.
  \label{eq:max_sum}
\end{equation}
More precisely, in the infinite-rank case the ratio is
\[
  O_p\!\left((\log n)^{-2}\right),
\]
while in the finite-rank case it is
\[
  O_p(n^{-1}).
\]
\end{proposition}


\begin{proof}

As in \Cref{lem:precision_concentration}, set
\[
  E_n:=(D_\nu+\eta_n I)^{-1/2}(D_{\nu,n}-D_\nu)(D_\nu+\eta_n I)^{-1/2}.
\]

As in \Cref{prop:regularized_gamma,prop:variance_consistency}, \Cref{lem:precision_concentration} provides us with the event $\{\norm{E_n}_{\op}\leq c<1\}$ which has probability tending to $1$ as $n\to\infty$. For the remainder of this proof, we work on this good event. 

By \Cref{lem:weight_operator_form},
\[
 \widehat\omega_i=\widehat t_{E_i}
 \inner{K_{X_i}}{(D_{\nu,n}+\eta_nI)^{-1}K_{x_0}}_{\F}.
\]
Applying Cauchy--Schwarz in the regularized inner product gives
\begin{align*}
 \widehat\omega_i^2
 &\leq\widehat t_{E_i}^2
 \norm{(D_{\nu,n}+\eta_n I)^{-1/2}K_{X_i}}_{\F}^2
 \norm{(D_{\nu,n}+\eta_n I)^{-1/2}K_{x_0}}_{\F}^2\\
 &=\widehat t_{E_i}^2\widehat\Ncal_{D,X_i}(\eta_n)
       \widehat\Ncal_{D,x_0}(\eta_n).
\end{align*}
On the good event, using Loewner order as in \Cref{prop:regularized_gamma},
\[
 D_{\nu,n}+\eta_nI\succeq(1-c)(D_\nu+\eta_n I).
\]
Moreover, \Cref{prop:noise_variances} implies that the two estimated precisions are bounded above and away from zero with probability tending to one. We now distinguish the two spectral regimes.

\medskip
\noindent\textbf{Infinite-rank case.}
The bounded-kernel leverage bound gives
\[
 \widehat\Ncal_{D,X_i}(\eta_n)\leq\kappa^2/\eta_n,
 \qquad
 \widehat\Ncal_{D,x_0}(\eta_n)
 \leq\frac1{1-c}\Ncal_{D,x_0}(\eta_n)=O(\eta_n^{-a}),
\]
where the last step is \Cref{lem:weighted_scale}. Therefore,
\[
 \max_i\widehat\omega_i^2=O_p(\eta_n^{-a-1}).
\]
On the other hand, since the estimated variances are bounded above and away from zero,
\[
 \sum_i\widehat\omega_i^2\asymp_p n\widehat V_n,
 \qquad \widehat V_n/V_n\to_p1,\qquad V_n\asymp\eta_n^{-a}.
\]
Consequently,
\[
 \frac{\max_i\widehat\omega_i^2}{\sum_i\widehat\omega_i^2}
 =O_p\!\left(\frac1{n\eta_n}\right)
 =O_p((\log n)^{-2})=o_p(1).
\]
Here $A_n\asymp_p B_n$ means that $A_n/B_n$ and $B_n/A_n$ are both bounded in probability.

\medskip
\noindent\textbf{Finite-rank case.}
Since $D_\nu\succeq(\min_gt_g)C_\nu$, its positive spectral gap is at least $(\min_gt_g)\mu_{\min}^+$. Thus $\Ncal_{D,X_i}(\eta_n)$ and $\Ncal_{D,x_0}(\eta_n)$ are bounded uniformly in $i$, and the same is true for their empirical versions on the good event. Hence $\max_i\widehat\omega_i^2=O_p(1)$. \Cref{prop:variance_consistency} gives $\widehat V_n\to_p V_0>0$, so $\sum_i\widehat\omega_i^2\asymp_p n$. Therefore the ratio is $O_p(n^{-1})=o_p(1)$.
\end{proof}



\newpage

% ============================================================
\begin{proposition}[Root-$m$ rate for the nuisance estimates]\label{prop:nuisance_rates}
\pdfbookmark[2]{Proposition \theproposition}{proposition.bookmark.\theproposition}
Suppose \Cref{assump:model} holds and $\tau_m=O(m^{-1/2})$. Then
\[
 \norm{\widehat c-c}+\max_g\{\norm{\widehat P_g-P_g}_{\op}
 +\norm{\widehat\mu_g-\mu_g}\}=O_p(m^{-1/2}).
\]
In particular, if $a_g=P_gc$, $b_g=-c^\top P_g\mu_g$, $\widehat a_g=\widehat P_g\widehat c$, and $\widehat b_g=-\widehat c^\top\widehat P_g\widehat\mu_g$, then
\[
 \max_g\{\norm{\widehat a_g-a_g}+|\widehat b_g-b_g|\}=O_p(m^{-1/2}).
\]
\end{proposition}
\begin{proof}
Gaussian fourth moments and fixed dimension give root-$m$ rates for the sample means and covariance matrices. Matrix inversion is continuously differentiable in a neighborhood of each positive definite $\Omega_g$, giving the precision rates. The first-fold mixture proportions have the same rate when they are estimated.

Collect these distribution parameters in a finite-dimensional vector $\alpha$. In a compact neighborhood of the truth, covariance eigenvalues stay bounded above and away from zero, and the mixture probabilities stay away from zero. The posterior environment probability is a logistic function of a quadratic polynomial in $x$. Its first derivatives therefore have polynomial envelopes in $\norm x$. The same is true of the derivatives of $V_\alpha(x,g)$. Consequently, the derivatives of $V_\alpha V_\alpha^\top$ and $V_\alpha Y$ have integrable envelopes, and the quantities at the true parameter have finite second moments. Here $f_0$ is bounded by $\kappa\norm{f_0}_{\F}$, and the Gaussian and sub-exponential variables have all polynomial moments.

The mean value theorem and empirical averages of these envelopes give
\begin{align*}
 \widehat M-M
 &=\frac1m\sum_{i\in I_1}(V_iV_i^\top-M)+O_p(\norm{\widehat\alpha-\alpha})
 =O_p(m^{-1/2}),\\
 \widehat h-\E[VY]
 &=\frac1m\sum_{i\in I_1}\{V_iY_i-\E[VY]\}
 +O_p(\norm{\widehat\alpha-\alpha})=O_p(m^{-1/2}).
\end{align*}
For fixed strata, center each summand at its within-environment expectation before summing; the weighted expectations give $M$ and $\E[VY]$. The same rate follows by a second-moment bound. No independence between $\widehat\alpha$ and these empirical moments is needed.

Since $M\succ0$, $(\widehat M+\tau_m I_d)^{-1}=O_p(1)$, and
\[
 \widehat c-c=(\widehat M+\tau_m I_d)^{-1}
 \{\widehat h-\E[VY]-(\widehat M-M)c-\tau_m c\}=O_p(m^{-1/2}).
\]
The rates for $a_g$ and $b_g$ follow by finite-dimensional multiplication.
\end{proof}

\begin{proposition}[Consistency of the environment variance estimates]\label{prop:noise_variances}
\pdfbookmark[2]{Proposition \theproposition}{proposition.bookmark.\theproposition}
Suppose \Cref{assump:model,assump:krr} hold and use the first-fold KRR and nuisance estimates described above. Then, for $g\in\{e,f\}$,
\[
 \widehat\sigma_{U,g}^2\to_p\sigma_{U,g}^2,\qquad \widehat t_g\to_p t_g.
\]
\end{proposition}
\begin{proof}
Write the first-fold corrected residual as $U_i+d_i$, where
\[
 d_i=f_0(X_i)-\widehat f_{\lambda_m}^{(1)}(X_i)
 +m_{E_i}(X_i)-\widehat m_{E_i}(X_i).
\]
By \Cref{thm:krr_rate}, boundedness of $K$, \Cref{prop:nuisance_rates}, and the empirical Gaussian second moments,
\[
 \frac1{m_g}\sum_{i\in I_1:E_i=g}d_i^2=o_p(1).
\]
The law of large numbers gives $m_g^{-1}\sum_{i\in I_1:E_i=g}U_i^2\to_p\sigma_{U,g}^2$. Cauchy--Schwarz makes the empirical cross term $o_p(1)$. Expanding $(U_i+d_i)^2$ proves variance consistency. Since $\sigma_{U,g}^2>0$, the lower truncation in \eqref{eq:sigma_hat} has no asymptotic effect, and inversion gives consistency of the precisions. This proof allows the same first-fold observations in all the fitted quantities.
\end{proof}

\newpage
\section{Additional Proofs}\label{sec:additional_proofs}

\begin{lemma}[Identification of the RKHS with the first power space]\label{lem:first_power_isometry}
\pdfbookmark[2]{Lemma \thelemma}{lemma.bookmark.\thelemma}
Under \Cref{assump:model}, the canonical map
\[
  I_\nu:\F\to[\F]^1_\nu,\qquad f\mapsto[f]_\nu,
\]
is an isometric isomorphism. In particular,
\[
  \norm{[f]_\nu}_{[\F]^1_\nu}=\norm{f}_{\F},\qquad f\in\F.
\]
\end{lemma}

\begin{proof}
By \Cref{assump:model}, $\nu$ has full support on $\X$ and $K$ is continuous. 

If there existed $0\neq f\in\ker I_\nu$, then $f=0$ $\nu$-a.e., but continuity of the kernel will give an open set $U\subset\X$ on which $|f|>0$. Full support then implies $\nu(U)>0$, a contradiction.

Hence $I_\nu$ is injective, so by the definition of the identifiable subspace,
\[
  \F=\F_\nu=(\ker I_\nu)^\perp=\overline{\ran(C_\nu)}.
\]
Therefore $(\mu_j^{1/2}e_j)_{j\in\mathcal J_\nu}$ is an orthonormal basis of $\F$. Every $f\in\F$ has a unique expansion
\[
  f=\sum_{j\in\mathcal J_\nu}b_j\mu_j^{1/2}e_j,
  \qquad
  (b_j)_{j\in\mathcal J_\nu}\in\ell_2(\mathcal J_\nu),
\]
where the sum is finite if $r_\nu<\infty$ and infinite if $r_\nu=\infty$. Moreover,
\[
  \norm{f}_{\F}^2=\sum_{j\in\mathcal J_\nu}b_j^2.
\]
Applying $I_\nu$ gives
\[
  [f]_\nu
  =
  \sum_{j\in\mathcal J_\nu}b_j\mu_j^{1/2}[e_j]_\nu,
\]
so the definition of the first power norm yields
\[
  \norm{[f]_\nu}_{[\F]^1_\nu}^2
  =
  \sum_{j\in\mathcal J_\nu}b_j^2
  =
  \norm{f}_{\F}^2.
\]
Conversely, every element of $[\F]^1_\nu$ has such an expansion, so $I_\nu:\F\to[\F]^1_\nu$ is onto and isometric.
\end{proof}

\phantomsection
\label{proof:gm_dual}
\pdfbookmark[2]{Proof of Proposition 4.1}{proof.bookmark.gm_dual}
\begin{proof}[Proof of \Cref{prop:gm_dual}]
We prove the two statements in order. We first derive the dual problem and the formula for the primal weights using the dual solution $\mu^\star$

\medskip


Expanding the RKHS alignment constraint using
\eqref{eq:alignment_gram_form} gives
\[
  \left\|
  K_{x_0}
  -\frac1n\sum_{i=1}^n\omega_iK_{X_i}
  \right\|_{\F}^2
  =
  \frac1{n^2}\omega^\top K\omega
  -\frac2n k_0^\top\omega
  +k_{00}.
\]
Hence the primal program \eqref{eq:gm_primal} can be written as
\[
  \min_{\omega\in\R^n}\frac1n\omega^\top W\omega
  \quad\text{subject to}\quad
  \frac1{n^2}\omega^\top K\omega
  -\frac2n k_0^\top\omega
  +k_{00}-\gamma^2
  \leq 0.
\]

\noindent
For a Lagrange multiplier $\mu\geq0$, the Lagrangian is
\[
  L(\omega,\mu)
  =
  \frac1n\omega^\top W\omega
  +
  \mu\left(
  \frac1{n^2}\omega^\top K\omega
  -\frac2n k_0^\top\omega
  +k_{00}-\gamma^2
  \right).
\]
For each fixed $\mu\geq0$, the map
$\omega\mapsto L(\omega,\mu)$ is strictly convex because
\[
  \frac1n W+\frac{\mu}{n^2}K
  \succ 0,
\]
since $K\succeq0$ and $n^{-1}W\succ0$. Therefore, the minimizer of
the Lagrangian with respect to $\omega$ is uniquely characterized by
the first-order condition.

\noindent Differentiating with respect to $\omega$ gives
\[
  \nabla_\omega L(\omega,\mu)
  =
  \frac2nW\omega
  +
  \frac{2\mu}{n^2}K\omega
  -
  \frac{2\mu}{n}k_0.
\]
Setting this equal to zero yields
\[
  \left(
  \frac1nW+\frac\mu{n^2}K
  \right)\omega
  =
  \frac\mu n k_0.
\]

\noindent which equals \eqref{eq:dual_weight_recovery}.

\noindent Since
\[
  \frac1nW+\frac\mu{n^2}K
\]
is positive definite and therefore invertible, the unique minimizer is
\begin{equation}
\label{eq:weight_formula}
  \omega^\star(\mu)
  =
  \frac\mu n
  \left(
  \frac1nW+\frac\mu{n^2}K
  \right)^{-1}
  k_0.
\end{equation}

\indent To obtain the dual objective function, first note that
the Lagrangian depending on $\omega$ can be written as
\[
  \omega^\top \left(\frac1nW+\frac\mu{n^2}K \right)\omega
  -
  \frac{2\mu}{n}k_0^\top\omega + \mu(k_{00}-\gamma^2).
\]
Using
\[
  \left(\frac1nW+\frac\mu{n^2}K \right)\omega^\star(\mu)
  =
  \frac\mu n k_0,
\]
we obtain
\[
  \omega^\star(\mu)^\top
  \left(\frac1nW+\frac\mu{n^2}K \right)
  \omega^\star(\mu)
  =
  \frac\mu n
  k_0^\top\omega^\star(\mu).
\]
Therefore, substituting in \eqref{eq:weight_formula}, we obtain 
\begin{align*}
  \inf_{\omega\in\R^n}L(\omega,\mu)
  &=
  \mu(k_{00}-\gamma^2)
  -
  \frac\mu n
  k_0^\top\omega^\star(\mu)
  \\
  &=
  \mu(k_{00}-\gamma^2)
  -
  \frac{\mu^2}{n^2}
  k_0^\top
  \left(
  \frac1nW+\frac\mu{n^2}K
  \right)^{-1}
  k_0.
\end{align*}
Hence the dual problem is
\[
  \max_{\mu\geq0}
  \left\{
  \mu(k_{00}-\gamma^2)
  -
  \frac{\mu^2}{n^2}
  k_0^\top
  \left(
  \frac1nW+\frac\mu{n^2}K
  \right)^{-1}
  k_0
  \right\},
\]
which is exactly \eqref{eq:gm_dual}.

By Slater's condition, strong duality holds because the primal problem
is convex and admits a strictly feasible point. Therefore, if
$\mu^\star$ denotes an optimal dual multiplier, the corresponding
primal optimizer is obtained by substituting $\mu^\star$ into
\eqref{eq:weight_formula}, giving
\[
  \widehat\omega
  =
  \frac{\mu^\star}{n}
  \left(
  \frac1nW+\frac{\mu^\star}{n^2}K
  \right)^{-1}
  k_0,
\]
which proves \eqref{eq:dual_weight_recovery}.

\medskip
\noindent
Secondly, we obtain an eigenbasis representation of our dual objective function and its derivatives.

Let
\[
 \widetilde K:=W^{-1/2}KW^{-1/2}=U\Lambda U^\top,
 \qquad a:=U^\top W^{-1/2}k_0.
\]
Then
\[
 \frac1nW+\frac\mu{n^2}K
 =W^{1/2}U\left(\frac1nI+\frac\mu{n^2}\Lambda\right)U^\top W^{1/2}.
\]
Consequently,
\[
 \frac{\mu^2}{n^2}k_0^\top
 \left(\frac1nW+\frac\mu{n^2}K\right)^{-1}k_0
 =\frac{\mu^2}{n^2}a^\top
 \left(\frac1nI+\frac\mu{n^2}\Lambda\right)^{-1}a
 =\sum_{j=1}^n\frac{a_j^2\mu^2}{n+\lambda_j\mu}.
\]
Thus the dual objective becomes
\[
  f(\mu)
  =
  (k_{00}-\gamma^2)\mu
  -
  \sum_{j=1}^n
  \frac{a_j^2\mu^2}{n+\lambda_j\mu},
\]
which is \eqref{eq:dual_eigenbasis}.

For each $j$,
\[
  \frac{d}{d\mu}
  \left(
  \frac{\mu^2}{n+\lambda_j\mu}
  \right)
  =
  \frac{
  2\mu(n+\lambda_j\mu)
  -
  \lambda_j\mu^2
  }{
  (n+\lambda_j\mu)^2
  }
  =
  \frac{
  \mu(2n+\lambda_j\mu)
  }{
  (n+\lambda_j\mu)^2
  }.
\]
Therefore,
\[
  f'(\mu)
  =
  k_{00}-\gamma^2
  -
  \sum_{j=1}^n
  a_j^2
  \frac{
  \mu(2n+\lambda_j\mu)
  }{
  (n+\lambda_j\mu)^2
  },
\]
which proves \eqref{eq:dual_first_derivative}.

Differentiating once more gives
\[
  \frac{d}{d\mu}
  \left[
  \frac{
  \mu(2n+\lambda_j\mu)
  }{
  (n+\lambda_j\mu)^2
  }
  \right]
  =
  \frac{2n^2}{(n+\lambda_j\mu)^3}.
\]
Hence
\[
  f''(\mu)
  =
  -
  \sum_{j=1}^n
  a_j^2
  \frac{2n^2}{(n+\lambda_j\mu)^3}
  \leq 0,
\]
which proves \eqref{eq:dual_second_derivative}. 
\end{proof}


\begin{lemma}[Regularized empirical-covariance concentration]\label{lem:operator_bernstein}
\pdfbookmark[2]{Lemma \thelemma}{lemma.bookmark.\thelemma}
Suppose Assumptions~\ref{assump:model} and \ref{assump:krr} hold. Define
\[
  L_\eta
  :=
  \begin{cases}
    \kappa^2/\eta, & r_\nu=\infty,\\[1mm]
    A_{\mathrm u}^{\mathrm{fin}}, & r_\nu<\infty,
  \end{cases}
\]
and
\[
  M_\eta
  :=
  \begin{cases}
    D\eta^{-p}, & r_\nu=\infty,\\[1mm]
    r_\nu, & r_\nu<\infty.
  \end{cases}
\]
Set
\[
  G_\eta
  :=\log\!\left(
    2eM_\eta\frac{\norm{C_\nu}_{\op}+\eta}{\norm{C_\nu}_{\op}}
  \right).
\]
For i.i.d. mixture observations, for every $\tau\ge1$, with probability at least $1-2e^{-\tau}$,
\begin{align}
  &\norm{(C_\nu+\eta I)^{-1/2}(C_\nu-C_{\nu,n})(C_\nu+\eta I)^{-1/2}}_{\op}
  \nonumber\\
  &\qquad\le
  \frac{4L_\eta\tau G_\eta}{3n}
  +\sqrt{
    \frac{2L_\eta\tau G_\eta}{n}
    \frac{\norm{C_\nu}_{\op}}{\norm{C_\nu}_{\op}+\eta}
  }.
  \label{eq:operator_bernstein}
\end{align}
For fixed strata, the same probability orders below follow by applying the inequality to each stratum and summing; only constants change. In particular:
\begin{enumerate}[label=(\roman*),leftmargin=2.2em]
  \item if $r_\nu=\infty$ and $\eta_n\asymp (\log n)^2/n$,
  \[
    \norm{(C_\nu+\eta_n I)^{-1/2}(C_\nu-C_{\nu,n})(C_\nu+\eta_n I)^{-1/2}}_{\op}
    =
    O_p\!\left(\sqrt{\frac{\log n}{n\eta_n}}\right)
    =O_p\!\left((\log n)^{-1/2}\right)
    =o_p(1);
  \]
  \item if $r_\nu<\infty$ and $\eta_n\asymp n^{-1}$,
  \[
    \norm{(C_\nu+\eta_n I)^{-1/2}(C_\nu-C_{\nu,n})(C_\nu+\eta_n I)^{-1/2}}_{\op}
    =
    O_p(n^{-1/2})
    =o_p(1).
  \]
\end{enumerate}
\end{lemma}





\begin{proof}
We adapt the self-adjoint operator Bernstein inequality stated in
\citet[Theorem~27]{fischer2020sobolev}. Their theorem is adapted from \cite{lincevher2020} and \cite{minsker2017extensions}.

Their theorem states the following.
Let $\xi$ be a random self-adjoint Hilbert--Schmidt operator satisfying
\[
    \norm{\xi}_{\op}\leq B
    \qquad\text{a.s.},
\]
and suppose that there exists a positive, self-adjoint, trace-class operator
$V$ such that
\[
    \E[\xi^2]\preceqL V.
\]
Then, defining
\[
    g(V)
    :=
    \log\!\left(
        2e\frac{\tr(V)}{\norm{V}_{\op}}
    \right),
\]
for every $\tau\geq1$, with probability at least $1-2e^{-\tau}$,
\begin{equation}
\label{eq:fischer_operator_bernstein}
    \left\|
        \frac1n\sum_{i=1}^n\xi_i-\E\xi
    \right\|_{\op}
    \leq
    \frac{4\tau B g(V)}{3n}
    +
    \sqrt{
        \frac{2\tau\norm{V}_{\op}g(V)}{n}
    }.
\end{equation}
We now identify the quantities $B$, $V$, and $g(V)$ in our setting.

For $x\in\X$, let
\[
    C_x:=K_x\otimes K_x
\]
and define the whitened random operator
\[
    \xi_i
    :=
    (C_\nu+\eta I)^{-1/2}
    C_{X_i}
    (C_\nu+\eta I)^{-1/2}.
\]
Writing
\[
    h_x:=(C_\nu+\eta I)^{-1/2}K_x,
\]
we have
\[
    \xi_i=h_{X_i}\otimes h_{X_i}.
\]
Thus $\xi_i$ is positive, self-adjoint, rank one, and hence
Hilbert--Schmidt, as required by Theorem~27.

We first identify the uniform operator bound $B$. Since
\[
    \norm{h\otimes h}_{\op}=\norm{h}_{\F}^2,
\]
we obtain
\begin{align*}
    \norm{\xi_i}_{\op}
    &=
    \norm{h_{X_i}}_{\F}^2 \\
    &=
    \left\langle
        K_{X_i},
        (C_\nu+\eta I)^{-1}K_{X_i}
    \right\rangle_{\F} \\
    &=
    \Ncal_{X_i}(\eta).
\end{align*}
By \eqref{eq:bounded_leverage} and \eqref{eq:uniform_leverage_finite}, in either spectral regime,
\[
    \Ncal_{X_i}(\eta)
    \leq
    L_\eta
    \qquad\text{a.s.}
\]
Therefore, we may take
\begin{equation}
\label{eq:bernstein_B_identification}
    B_\eta
    :=
    L_\eta.
\end{equation}

Next, the expectation of $\xi_i$ is
\begin{align*}
    \E[\xi_i]
    &=
    (C_\nu+\eta I)^{-1/2}
    \E[C_{X_i}]
    (C_\nu+\eta I)^{-1/2} \\
    &=
    (C_\nu+\eta I)^{-1/2}
    C_\nu
    (C_\nu+\eta I)^{-1/2}.
\end{align*}
Since $\xi_i=h_{X_i}\otimes h_{X_i}$ is rank one,
\[
    \xi_i^2
    =
    \norm{h_{X_i}}_{\F}^2\xi_i.
\]
Using the bound $\Ncal_{X_i}(\eta)\le L_\eta$,
\[
    \xi_i^2
    \preceqL
    L_\eta\xi_i.
\]
Taking expectations gives
\[
    \E[\xi_i^2]
    \preceqL
    L_\eta\E[\xi_i].
\]
Hence the operator $V$ can be chosen as
\begin{equation}
\label{eq:bernstein_V_identification}
    V_\eta
    :=
    L_\eta\E[\xi_i].
\end{equation}
Since $\E[\xi_i]$ is positive, self-adjoint, and trace class, so is
$V_\eta$.

It remains to identify $\norm{V_\eta}_{\op}$ and $g(V_\eta)$. We have
\[
    \norm{\E[\xi_i]}_{\op}
    =
    \sup_j\frac{\mu_j}{\mu_j+\eta}
    =
    \frac{\norm{C_\nu}_{\op}}
         {\norm{C_\nu}_{\op}+\eta}.
\]
Therefore,
\begin{equation}
\label{eq:bernstein_V_norm}
    \norm{V_\eta}_{\op}
    =
    L_\eta
    \frac{\norm{C_\nu}_{\op}}
         {\norm{C_\nu}_{\op}+\eta}.
\end{equation}
Moreover,
\[
    \tr(\E[\xi_i])
    =
    \tr\!\left(
        C_\nu(C_\nu+\eta I)^{-1}
    \right)
    =
    \Ncal_\nu(\eta)
    \leq
    M_\eta,
\]
where the final inequality follows from \Cref{assump:krr}: in the infinite-rank case $M_\eta=D\eta^{-p}$, while in the finite-rank case $\Ncal_\nu(\eta)\le r_\nu=M_\eta$. We finally obtain
\begin{align*}
    g(V_\eta)
    &=
    \log\!\left(
        2e
        \frac{\tr(V_\eta)}
             {\norm{V_\eta}}
    \right) \\
    &=
    \log\!\left(
        2e
        \frac{\tr(\E[\xi_i])}
             {\norm{\E[\xi_i]}_{\op}}
    \right) \\
    &\leq
    \log\!\left(
        2eM_\eta
        \frac{\norm{C_\nu}_{\op}+\eta}
             {\norm{C_\nu}_{\op}}
    \right)
    =
    G_\eta.
\end{align*}

Finally, observe that
\begin{align*}
    \frac1n\sum_{i=1}^n\xi_i- \E[\xi_i]
    &=
    (C_\nu+\eta I)^{-1/2}
    \left(
        \frac1n\sum_{i=1}^n C_{X_i}-C_\nu
    \right)
    (C_\nu+\eta I)^{-1/2} \\
    &=
    (C_\nu+\eta I)^{-1/2}
    (C_{\nu,n}-C_\nu)
    (C_\nu+\eta I)^{-1/2}.
\end{align*}

We therefore have
\[
    B=B_\eta
    =
    L_\eta,
    \qquad
    V=V_\eta^{},
    \qquad
    g(V_\eta^{})\leq G_\eta.
\]
Substituting \eqref{eq:bernstein_B_identification} and
\eqref{eq:bernstein_V_norm} into
\eqref{eq:fischer_operator_bernstein} yields, with probability at least
$1-2e^{-\tau}$,
\begin{align*}
  &\norm{
    (C_\nu+\eta I)^{-1/2}
    (C_{\nu,n} - C_\nu)
    (C_\nu+\eta I)^{-1/2}
  }_{\op}
  \\
  &\qquad\leq
  \frac{4L_\eta\tau G_\eta}{3n}
  +
  \sqrt{
    \frac{2L_\eta\tau G_\eta}{n}
    \frac{\norm{C_\nu}_{\op}}
         {\norm{C_\nu}_{\op}+\eta}
  },
\end{align*}

which is \eqref{eq:operator_bernstein}.

For the final assertion, let $\eta_n\asymp (\log n)^2/n$ if $r_\nu=\infty$ and $\eta_n\asymp n^{-1}$ if $r_\nu<\infty$.

\medskip
\noindent\textbf{Infinite-rank case.}
Then
\[
  L_{\eta_n}=\frac{\kappa^2}{\eta_n}
  \asymp \frac{n}{(\log n)^2},
  \qquad
  M_{\eta_n}=D\eta_n^{-p}\asymp \frac{n^p}{(\log n)^{2p}},
\]
and
\[
  \frac{\norm{C_\nu}_{\op}+\eta_n}{\norm{C_\nu}_{\op}}
  =O(1),
\]
so
\[
  G_{\eta_n}=O(\log n).
\]
The first term in \eqref{eq:operator_bernstein} is therefore
\[
  O_p\!\left(\frac{\log n}{n\eta_n}\right)
  =O_p\!\left((\log n)^{-1}\right),
\]
while the second term is
\[
  O_p\!\left(
  \sqrt{\frac{\log n}{n\eta_n}}
  \right)
  =O_p\!\left((\log n)^{-1/2}\right).
\]
The second term dominates asymptotically, and hence
\[
  \norm{
    (C_\nu+\eta_n I)^{-1/2}
    (C_\nu-C_{\nu,n})
    (C_\nu+\eta_n I)^{-1/2}
  }_{\op}
  =
  O_p\!\left(
    \sqrt{\frac{\log n}{n\eta_n}}
  \right)
  =O_p\!\left((\log n)^{-1/2}\right)
  =o_p(1).
\]

\medskip
\noindent\textbf{Finite-rank case.}
Here
\[
  L_{\eta_n}=A_{\mathrm u}^{\mathrm{fin}}=O(1),
  \qquad
  M_{\eta_n}=r_\nu=O(1),
\]
and therefore $G_{\eta_n}=O(1)$. Consequently,
\[
  \norm{
    (C_\nu+\eta_n I)^{-1/2}
    (C_\nu-C_{\nu,n})
    (C_\nu+\eta_n I)^{-1/2}
  }_{\op}
  =
  O_p(n^{-1/2})
  =o_p(1).
\]
\end{proof}




For completeness, in the stratified case apply the same Bernstein argument within each environment after whitening by $(C_\nu+\eta I)^{-1/2}$. Since $C_g\preceq\pi_g^{-1}C_\nu$, the whitened expectation has bounded operator norm. Its trace is at most $\pi_g^{-1}\Ncal_\nu(\eta)$. Thus the envelopes and logarithmic factors have the same orders, with constants depending on the fixed positive $\pi_g$. Summing the two centered errors proves the stated orders for $C_{\nu,n}-C_\nu$.

\begin{lemma}[Operator form of the debiasing weights]
\pdfbookmark[2]{Lemma \thelemma}{lemma.bookmark.\thelemma}
\label{lem:weight_operator_form}
Let $S_n$ be the sampling operator from \Cref{defn:sampling_operator}. 
If the optimal dual multiplier is $\mu^\star=\eta^{-1}$, then
\begin{equation}
 \widehat\omega_i=\widehat t_{E_i}
 \left((D_{\nu,n}+\eta I)^{-1}K_{x_0}\right)(X_i),\qquad i=1,\ldots,n.
 \label{eq:weight_operator_form}
\end{equation}
Equivalently,
\begin{equation}
 \widehat\omega=n(K+n\eta W)^{-1}k_0.
 \label{eq:weight_matrix_form}
\end{equation}
\end{lemma}

\begin{proof}
We first establish the intertwining identity for the bounded operators $\widetilde S_n^*\widetilde S_n$ and $\widetilde S_n\widetilde S_n^*$:
\begin{equation}
  \widetilde S_n h(\widetilde S_n^*\widetilde S_n)
  =
  h(\widetilde S_n\widetilde S_n^*)\widetilde S_n
  \label{eq:intertwining}
\end{equation}
for every continuous function $h$ on a compact interval containing the spectra of
$\widetilde S_n^*\widetilde S_n$ and $\widetilde S_n\widetilde S_n^*$.

For the monomial $h(t)=t^m$, we claim that
\[
  \widetilde S_n(\widetilde S_n^*\widetilde S_n)^m
  =
  (\widetilde S_n\widetilde S_n^*)^m\widetilde S_n,
  \qquad m\ge0.
\]
The case $m=0$ is immediate. If the identity holds for some $m$, then
\begin{align*}
  \widetilde S_n(\widetilde S_n^*\widetilde S_n)^{m+1}
  &=
  \widetilde S_n(\widetilde S_n^*\widetilde S_n)^m(\widetilde S_n^*\widetilde S_n) \\
  &=
  (\widetilde S_n\widetilde S_n^*)^m \widetilde S_n\widetilde S_n^*\widetilde S_n \\
  &=
  (\widetilde S_n\widetilde S_n^*)^{m+1}\widetilde S_n.
\end{align*}
Thus the identity holds for all $m$ by induction. By linearity of the functional
calculus, it therefore holds for every polynomial $p$:
\[
  \widetilde S_np(\widetilde S_n^*\widetilde S_n)
  =
  p(\widetilde S_n\widetilde S_n^*)\widetilde S_n.
\]

Now let $h$ be continuous on a compact interval $[0,M]$ containing both spectra.
By the Weierstrass approximation theorem, there exist polynomials $p_j$ such that
\[
  \sup_{t\in[0,M]}|p_j(t)-h(t)|\to0.
\]
Since $\widetilde S_n^*\widetilde S_n$ and $\widetilde S_n\widetilde S_n^*$ are bounded, this implies
\[
  \widetilde S_np_j(\widetilde S_n^*\widetilde S_n)\to \widetilde S_nh(\widetilde S_n^*\widetilde S_n),
  \qquad
  p_j(\widetilde S_n\widetilde S_n^*)\widetilde S_n\to h(\widetilde S_n\widetilde S_n^*)\widetilde S_n
\]
in operator norm. Taking limits in
\[
  \widetilde S_np_j(\widetilde S_n^*\widetilde S_n)
  =
  p_j(\widetilde S_n\widetilde S_n^*)\widetilde S_n
\]
thus proves \eqref{eq:intertwining}.

For $\eta > 0$, take now the function
\[
  h(t)=\frac{1}{t+\eta}.
\]
Since
\[
 \widetilde S_n^*\widetilde S_n=D_{\nu,n},\qquad
 \widetilde S_n\widetilde S_n^*=\frac1n\widetilde K,\qquad
 \widetilde S_nK_{x_0}=W^{-1/2}k_0,
\]
the intertwining identity yields
\begin{align*}
 W^{-1/2}\widetilde S_n(D_{\nu,n}+\eta I)^{-1}K_{x_0}
 &=nW^{-1/2}(\widetilde K+n\eta I)^{-1}W^{-1/2}k_0\\
 &=n(K+n\eta W)^{-1}k_0.
\end{align*}
The left-hand side is $W^{-1}S_n(D_{\nu,n}+\eta I)^{-1}K_{x_0}$, whose $i$th component is the expression in \eqref{eq:weight_operator_form}. Finally, substituting $\mu^\star=\eta^{-1}$ into the dual recovery formula gives the same matrix expression.
\end{proof}



\begin{lemma}[Expectation of truncated second moment]
\pdfbookmark[2]{Lemma \thelemma}{lemma.bookmark.\thelemma}
\label{lem:truncated_second_moment}
Suppose that the errors $U_i$ satisfy
\[
\E[U_i\mid X_i,E_i]=0,
\qquad
\E[|U_i|^j\mid X_i,E_i]
\leq
\frac{j!}{2}B_U^2 L^{j-2},
\qquad j\geq 2.
\]
Then, for every $t\geq 0$,
\begin{equation}
\label{eq:truncated_2nd_moment_tail_bound}
\E\!\left[
    U_i^2\bm{1}\{|U_i|>t\}
    \,\middle|\,X_i,E_i
\right]
\leq
C(B_U^2+L^2)
\exp\!\left\{
    -c\min\left(
        \frac{t^2}{B_U^2},
        \frac{t}{L}
    \right)
\right\},
\end{equation}
where the constants $C,c>0$ are independent of $t$.
\end{lemma}

\begin{proof}
From the Bernstein condition, for $0<\lambda<1/L$,
\begin{align*}
\E[e^{\lambda U_i}\mid X_i,E_i]
&=
1+\sum_{j=2}^\infty
\frac{\lambda^j}{j!}\E[U_i^j\mid X_i,E_i]
\\
&\leq
1+\frac{1}{2}
\sum_{j=2}^\infty
\lambda^jB_U^2L^{j-2}
\\
&=
1+\frac{B_U^2\lambda^2}{2}
\sum_{j=2}^\infty(\lambda L)^{j-2}
\\
&=
1+
\frac{B_U^2\lambda^2}
     {2(1-\lambda L)}
\\
&\leq
\exp\!\left\{
    \frac{B_U^2\lambda^2}
         {2(1-\lambda L)}
\right\}.
\end{align*}
The interchange between expectation and summation is justified by
absolute convergence under the Bernstein moment condition. The last
inequality follows from $1+x\leq e^x$.

By Chernoff's inequality,
\[
\Pp(U_i>t\mid X_i,E_i)
\leq
\exp\!\left\{
    -\lambda t+
    \frac{B_U^2\lambda^2}
         {2(1-\lambda L)}
\right\}.
\]
Setting
\[
\lambda
=
\frac{t}{B_U^2+Lt},
\]
we obtain
\[
\Pp(U_i>t\mid X_i,E_i)
\leq
\exp\!\left\{
    -\frac{t^2}
           {2(B_U^2+Lt)}
\right\}.
\]
Applying the same argument to $-U_i$ gives the two-sided
sub-exponential tail bound
\begin{equation}
\label{eq:bernstein_error_tail}
\Pp(|U_i|>t\mid X_i,E_i)
\leq
2\exp\!\left\{
    -c\min\left(
        \frac{t^2}{B_U^2},
        \frac{t}{L}
    \right)
\right\},
\end{equation}
for some constant $c>0$.

From the tail-integration formula,
\begin{align*}
\E\!\left[
    U_i^2\bm{1}\{|U_i|>t\}
    \,\middle|\,X_i,E_i
\right]
&=
t^2\Pp(|U_i|>t\mid X_i,E_i)
+
2\int_t^\infty
u\,\Pp(|U_i|>u\mid X_i,E_i)\,du
\\
&\leq
\underbrace{
2t^2
\exp\!\left\{
    -c\min\left(
        \frac{t^2}{B_U^2},
        \frac{t}{L}
    \right)
\right\}
}_{(A)}
\\
&\qquad
+
\underbrace{
4\int_t^\infty
u
\exp\!\left\{
    -c\min\left(
        \frac{u^2}{B_U^2},
        \frac{u}{L}
    \right)
\right\}
\,du
}_{(B)}.
\end{align*}

For $(A)$, consider separately the cases
\[
0\leq t\leq \frac{B_U^2}{L}
\qquad\text{and}\qquad
t>\frac{B_U^2}{L}.
\]
In the first case,
\[
A
=
2t^2e^{-ct^2/B_U^2}
\leq
C_1B_U^2
e^{-c_1t^2/B_U^2},
\]
while in the second case,
\[
A
=
2t^2e^{-ct/L}
\leq
C_1L^2e^{-c_1t/L},
\]
where $C_1,c_1>0$ are independent of $t$. Hence,
\[
A
\leq
C_1(B_U^2+L^2)
\exp\!\left\{
    -c_1\min\left(
        \frac{t^2}{B_U^2},
        \frac{t}{L}
    \right)
\right\}.
\]

Similarly, for $(B)$, consider the same two cases. If
$t\leq B_U^2/L$, split the integral at
$B_U^2/L$:
\begin{align*}
B
&\leq
4\int_t^{B_U^2/L}
u e^{-cu^2/B_U^2}\,du
+
4\int_{B_U^2/L}^\infty
u e^{-cu/L}\,du
\\
&\leq
C_2(B_U^2+L^2)
e^{-c_2t^2/B_U^2}.
\end{align*}
If $t>B_U^2/L$, then
\begin{align*}
B
&\leq
4\int_t^\infty
u e^{-cu/L}\,du
\\
&\leq
C_2L^2e^{-c_2t/L}.
\end{align*}
Therefore,
\[
B
\leq
C_2(B_U^2+L^2)
\exp\!\left\{
    -c_2\min\left(
        \frac{t^2}{B_U^2},
        \frac{t}{L}
    \right)
\right\}.
\]

Combining the bounds for $(A)$ and $(B)$ and adjusting the constants
$C,c>0$ gives
\[
\E\!\left[
    U_i^2\bm{1}\{|U_i|>t\}
    \,\middle|\,X_i,E_i
\right]
\leq
C(B_U^2+L^2)
\exp\!\left\{
    -c\min\left(
        \frac{t^2}{B_U^2},
        \frac{t}{L}
    \right)
\right\},
\]
which proves the claim.
\end{proof}



% ============================================================

\begin{lemma}[Spectral and source conditions for the mixture]\label{lem:mixture_source}
\pdfbookmark[2]{Lemma \thelemma}{lemma.bookmark.\thelemma}
Suppose the two-sided eigenvalue condition and the source condition in \Cref{assump:krr} hold separately under $\nu_g$, with exponents $p_g$ and $\beta_g$. Then they hold for $\nu$ with $p=\max_gp_g$ and $\beta=\min_g\beta_g$.
\end{lemma}
\begin{proof}
Since $C_\nu\succeq\pi_gC_g$, the min--max principle gives the lower eigenvalue bound using an environment with exponent $p$. For the upper bound, the intersection of two subspaces of codimension $j-1$ has codimension at most $2j-2$, giving
\[
 \mu_{2j-1}(C_\nu)\leq\pi_e\mu_j(C_e)+\pi_f\mu_j(C_f)\lesssim j^{-1/p}.
\]
Monotonicity gives the upper bound for all indices.

Put $s=(\beta-1)/2\in(0,1/2]$. The first-power isometry and the environmental source conditions give $f_0=C_g^sv_g$ for fixed $v_g\in\F$. For $0<r<1$, the scalar identity
\[
 t^r=\frac{\sin(\pi r)}\pi\int_0^\infty\frac{t}{t+u}u^{r-1}\,du
\]
and inverse-order monotonicity show that positive operator powers preserve Loewner order for $0<r\leq1$. Apply this with $r=2s$ to get
\[
 C_g^{2s}\preceq\pi_g^{-2s}C_\nu^{2s},\qquad
 \norm{C_g^sz}_{\F}\leq\pi_g^{-s}\norm{C_\nu^sz}_{\F}.
\]
The map $C_\nu^sz\mapsto C_g^sz$ extends to a bounded map on the closure of its range. Taking its adjoint gives $C_g^s=C_\nu^sB_g$ for a bounded operator $B_g$. Hence $f_0=C_\nu^sB_gv_g$, which is the pooled source condition. No commutation of the environment operators is required.
\end{proof}

\begin{lemma}[Regularized precision-weighted covariance concentration]\label{lem:precision_concentration}
\pdfbookmark[2]{Lemma \thelemma}{lemma.bookmark.\thelemma}
Suppose \Cref{assump:model,assump:krr} hold, and construct the variance estimates as in \eqref{eq:sigma_hat}. If $\eta_n\asymp(\log n)^2/n$ in the infinite-rank case or $\eta_n\asymp n^{-1}$ in the finite-rank case, then
\[
 E_n:=(D_\nu+\eta_n I)^{-1/2}(D_{\nu,n}-D_\nu)(D_\nu+\eta_n I)^{-1/2},
 \qquad \norm{E_n}_{\op}=o_p(1).
\]
Moreover,
\begin{align}
 &\norm{(D_\nu+\eta_n I)^{1/2}
 \{(D_{\nu,n}+\eta_n I)^{-1}-(D_\nu+\eta_n I)^{-1}\}K_{x_0}}_{\F}\nonumber\\
 &\qquad=o_p\!\left(\sqrt{\Ncal_{D,x_0}(\eta_n)}\right).
 \label{eq:representer_perturbation}
\end{align}
\end{lemma}
\begin{proof}
First use the true precisions in the empirical operator. Let $l=\min_gt_g>0$ and $u=\max_gt_g<\infty$. The whitened summands are
\[
 t_{E_i}(D_\nu+\eta I)^{-1/2}(K_{X_i}\otimes K_{X_i})(D_\nu+\eta I)^{-1/2}.
\]
Their norms are at most $u\kappa^2/\eta$, and their expected operator has norm at most one. Its trace is at most $\operatorname{tr}(D_\nu)/\eta$, giving a logarithmic factor $O(\log(1/\eta))$. The same operator Bernstein calculation as in \Cref{lem:operator_bernstein} gives
\[
 O_p\!\left(\sqrt{\frac{\log n}{n\eta_n}}+
 \frac{\log n}{n\eta_n}\right)=o_p(1)
\]
in the infinite-rank case. In the finite-rank case, the positive spectral gap gives $O_p(n^{-1/2})$. For fixed strata, apply the bound separately within each environment and sum.

By \Cref{prop:noise_variances}, $\max_g|\widehat t_g/t_g-1|=o_p(1)$. The whitened error from substituting estimated precisions is bounded in norm by this maximum times the norm of the whitened empirical operator with true precisions, which is $O_p(1)$. This proves $\norm{E_n}_{\op}=o_p(1)$.

Finally,
\[
 (D_\nu+\eta I)^{1/2}(D_{\nu,n}+\eta I)^{-1}K_{x_0}
 =(I+E_n)^{-1}(D_\nu+\eta I)^{-1/2}K_{x_0}.
\]
On $\norm{E_n}_{\op}<1$, subtracting the population expression bounds the norm in \eqref{eq:representer_perturbation} by
\[
 \frac{\norm{E_n}_{\op}}{1-\norm{E_n}_{\op}}
 \sqrt{\Ncal_{D,x_0}(\eta_n)},
\]
which proves the claim.
\end{proof}

\begin{lemma}[Variance scale under the pooled leverage assumption]\label{lem:weighted_scale}
\pdfbookmark[2]{Lemma \thelemma}{lemma.bookmark.\thelemma}
Suppose \Cref{assump:model,assump:local_leverage} hold and $r_\nu=\infty$. Then
\[
 \Ncal_{D,x_0}(\eta)\asymp\Ncal_{x_0}(\eta)\asymp\eta^{-a},
 \qquad
 \inner{(D_\nu+\eta I)^{-1}K_{x_0}}
 {D_\nu(D_\nu+\eta I)^{-1}K_{x_0}}_{\F}\asymp\eta^{-a}.
\]
\end{lemma}
\begin{proof}
Let $l=\min_gt_g>0$ and $u=\max_gt_g<\infty$. Since $lC_\nu\preceq D_\nu\preceq uC_\nu$, inversion gives
\[
 \frac1u\Ncal_{x_0}(\eta/u)\leq\Ncal_{D,x_0}(\eta)
 \leq\frac1l\Ncal_{x_0}(\eta/l).
\]
Thus, for some constants $L_0,U_0>0$ and all small $\eta$,
$L_0\eta^{-a}\leq\Ncal_{D,x_0}(\eta)\leq U_0\eta^{-a}$.
Let $b_j:=\inner{K_{x_0}}{\psi_j}_{\F}$. Write the variance quadratic form as
\[
 V(\eta)=\sum_j\frac{\rho_jb_j^2}{(\rho_j+\eta)^2}\leq\Ncal_{D,x_0}(\eta).
\]
Next, we use a helper function to sandwich our target, as in the unweighted variance argument. For fixed $b>1$,
\begin{align*}
 b\Ncal_{D,x_0}(b\eta)-\Ncal_{D,x_0}(\eta)
 &=(b-1)\sum_j\frac{\rho_jb_j^2}{(\rho_j+\eta)(\rho_j+b\eta)}\\
 &\leq(b-1)V(\eta).
\end{align*}
Choose $b$ large enough that $L_0b^{1-a}>U_0$. Then
\[
 V(\eta)\geq\frac{L_0b^{1-a}-U_0}{b-1}\eta^{-a}.
\]
This proves the variance comparison using only regular variation for $C_\nu$. It also implies $\eta\Ncal_{D,x_0}(\eta)\to0$.
\end{proof}

\begin{lemma}[Vanishing normalized population weighted moments]\label{lem:weak_representer}
\pdfbookmark[2]{Lemma \thelemma}{lemma.bookmark.\thelemma}
Suppose \Cref{assump:model,assump:local_leverage} hold and $r_\nu=\infty$. For every fixed $b=b(X,E)\in L_2(\mathbb P_{X,E})$,
\[
 \frac{\E[t_E\{(D_\nu+\eta I)^{-1}K_{x_0}\}(X)b(X,E)]}{\sqrt{V(\eta)}}\longrightarrow0,
\]
where $V(\eta)$ is the population variance quadratic form in \Cref{lem:weighted_scale}.
\end{lemma}
\begin{proof}
We use the $L_2$ space for the finite measure $t_E\,d\mathbb P_{X,E}$ and write its inner product as $\inner{\cdot}{\cdot}_t$. The bounds $0<l\leq t_E\leq u<\infty$ make its norm equivalent to the original $L_2$ norm. Let $\mathcal H_t$ be the closure of the functions $f(X)$ with $f\in\F$ in this space. The functions
\[
 z_\eta(X,E):=\frac{\{(D_\nu+\eta I)^{-1}K_{x_0}\}(X)}{\sqrt{V(\eta)}}
\]
belong to $\mathcal H_t$ and have norm one. For every fixed $f\in\F$,
\begin{align*}
 \inner{z_\eta}{f}_t
 &=\frac{\inner{(D_\nu+\eta I)^{-1}K_{x_0}}{D_\nu f}_{\F}}{\sqrt{V(\eta)}}\\
 &=\frac{f(x_0)-\eta\inner{(D_\nu+\eta I)^{-1}K_{x_0}}{f}_{\F}}{\sqrt{V(\eta)}}.
\end{align*}
The numerator is bounded, since
\[
 \eta\norm{(D_\nu+\eta I)^{-1}K_{x_0}}_{\F}
 \leq\sqrt{\eta\Ncal_{D,x_0}(\eta)}\longrightarrow0.
\]
The denominator diverges by \Cref{lem:weighted_scale}. Thus these inner products tend to zero. Density and the unit-norm bound extend this to all of $\mathcal H_t$. For arbitrary $b$ in the weighted $L_2$ space, its component orthogonal to $\mathcal H_t$ has zero inner product with $z_\eta$; projection proves the claim. In particular, it is not necessary that affine functions belong to $\F$, or that $\F$ be dense in all of $L_2(\nu)$.
\end{proof}

\begin{lemma}[Vanishing normalized empirical weighted moments]\label{lem:weighted_moments}
\pdfbookmark[2]{Lemma \thelemma}{lemma.bookmark.\thelemma}
Suppose Assumptions~\ref{assump:model}--\ref{assump:local_leverage} hold, $r_\nu=\infty$, and use the weights from \Cref{prop:regularized_gamma}. For every fixed $b=b(X,E)\in L_2(\mathbb P_{X,E})$,
\[
 \frac1n\sum_{i=1}^n\widehat\omega_i b(X_i,E_i)=o_p(\sqrt{V_n}).
\]
\end{lemma}
\begin{proof}
First use the population weights
\[
 \omega_i^\circ:=t_{E_i}\{(D_\nu+\eta_n I)^{-1}K_{x_0}\}(X_i).
\]
By regularized Cauchy--Schwarz and boundedness of the kernel,
\[
 |\{(D_\nu+\eta I)^{-1}K_{x_0}\}(x)|^2
 \leq\frac{\kappa^2}{\eta}\Ncal_{D,x_0}(\eta).
\]
Independence of the second-fold observations gives
\[
 \operatorname{Var}\!\left(\frac1n\sum_i\omega_i^\circ b(X_i,E_i)\right)
 \leq\frac{u^2\kappa^2\Ncal_{D,x_0}(\eta_n)}{n\eta_n}\E[b(X,E)^2]
 =o(V_n).
\]
The same bound holds for the fixed-stratum mean centered at its pooled expectation. \Cref{lem:weak_representer} makes that expectation $o(\sqrt{V_n})$, proving the result for population weights.

It remains to replace them by empirical weights. Write
\begin{align*}
 \widehat\omega_i-\omega_i^\circ
 &=\widehat t_{E_i}\big\{(D_{\nu,n}+\eta_n I)^{-1}K_{x_0}
              -(D_\nu+\eta_n I)^{-1}K_{x_0}\big\}(X_i)\\
 &\quad+(\widehat t_{E_i}-t_{E_i})\{(D_\nu+\eta_n I)^{-1}K_{x_0}\}(X_i).
\end{align*}
For the first term, \Cref{lem:precision_concentration} and
$D_{\nu,n}\preceq(1+\norm{E_n}_{\op})(D_\nu+\eta_n I)$ give an empirical squared norm of $o_p(\Ncal_{D,x_0}(\eta_n))$. For the second, factor out $\max_g|\widehat t_g-t_g|^2=o_p(1)$; the remaining empirical squared norm is $O_p(\Ncal_{D,x_0}(\eta_n))$ by the same operator comparison. The estimated precisions are bounded above and away from zero with probability tending to one. Consequently,
\[
 \frac1n\sum_i(\widehat\omega_i-\omega_i^\circ)^2
 =o_p(\Ncal_{D,x_0}(\eta_n))=o_p(V_n).
\]
Since $n^{-1}\sum_i b(X_i,E_i)^2=O_p(1)$, empirical Cauchy--Schwarz proves the claim.
\end{proof}

\addcontentsline{toc}{section}{References}
% ============================================================

\begingroup
\setstretch{1.1}
\begin{thebibliography}{99}

\bibitem[Lin and Cevher(2020)]{lincevher2020}
J. Lin and V. Cevher.
\newblock Optimal convergence for distributed learning with stochastic gradient methods and spectral algorithms.
\newblock \emph{Journal of Machine Learning Research}, 21(147):1--63, 2020.

\bibitem[Blanchard and M\"ucke(2018)]{blanchard2018optimal}
G. Blanchard and N. M\"ucke.
\newblock Optimal rates for regularization of statistical inverse learning problems.
\newblock \emph{Foundations of Computational Mathematics}, 18(4):971--1013, 2018.

\bibitem[Caponnetto and De Vito(2007)]{caponnetto2007optimal}
A. Caponnetto and E. De Vito.
\newblock Optimal rates for the regularized least-squares algorithm.
\newblock \emph{Foundations of Computational Mathematics}, 7:331--368, 2007.

\bibitem[Fischer and Steinwart(2020)]{fischer2020sobolev}
S. Fischer and I. Steinwart.
\newblock Sobolev norm learning rates for regularized least-squares algorithms.
\newblock \emph{Journal of Machine Learning Research}, 21(205):1--38, 2020.

\bibitem[Lin et~al.(2020)Lin, Rudi, Rosasco, and Cevher]{lin2020optimal}
J. Lin, A. Rudi, L. Rosasco, and V. Cevher.
\newblock Optimal rates for spectral algorithms with least-squares regression over Hilbert spaces.
\newblock \emph{Applied and Computational Harmonic Analysis}, 48(3):868--890, 2020.

\bibitem[Minsker(2017)]{minsker2017extensions}
S. Minsker.
\newblock On some extensions of Bernstein's inequality for self-adjoint operators.
\newblock \emph{Statistics \& Probability Letters}, 127:111--119, 2017.

\bibitem[Pauwels et~al.(2018)Pauwels, Bach, and Vert]{pauwels2018relating}
E. Pauwels, F. Bach, and J.-P. Vert.
\newblock Relating leverage scores and density using regularized Christoffel functions.
\newblock In \emph{Advances in Neural Information Processing Systems}, volume 31, 2018.
\bibitem[Wahba(1990)]{wahba1990spline}
G.~Wahba.
\newblock \emph{Spline Models for Observational Data}.
\newblock CBMS-NSF Regional Conference Series in Applied Mathematics, Vol.~59.
\newblock SIAM, Philadelphia, 1990.

\bibitem[Kimeldorf and Wahba(1971)]{kimeldorf1971some}
G.~S.~Kimeldorf and G.~Wahba.
\newblock Some results on Tchebycheffian spline functions.
\newblock \emph{Journal of Mathematical Analysis and Applications},
33(1):82--95, 1971.

\end{thebibliography}
\endgroup

\end{document}